\documentclass[aps,prd,twocolumn,superscriptaddress,nofootinbib]{revtex4-2}

\usepackage{amsmath}
\usepackage{mathtools}
\usepackage{amssymb}
\usepackage{amsthm}
\usepackage{booktabs}
\usepackage{array}
\usepackage{xcolor}
\usepackage{graphicx}
\usepackage{mathrsfs}

\emergencystretch=2em

\newtheorem{theorem}{Theorem}[section]
\newtheorem{proposition}[theorem]{Proposition}
\newtheorem{lemma}[theorem]{Lemma}
\newtheorem{corollary}[theorem]{Corollary}
\newtheorem{conjecture}[theorem]{Conjecture}
\newtheorem{axiom}[theorem]{Axiom}
\theoremstyle{definition}
\newtheorem{definition}[theorem]{Definition}
\theoremstyle{remark}
\newtheorem{remark}[theorem]{Remark}
\newtheorem{example}[theorem]{Example}

\usepackage[hyperfootnotes=false]{hyperref}
\pdfstringdefDisableCommands{%
    \def\\{ }%
}
\hypersetup{
    pdftitle={Arithmetic Positivity and Absolute Hodge Classes: Why Hodge-Riemann Positivity Does Not Strengthen Deligne's Absoluteness},
    pdfauthor={Lukas Geiger},
    pdfsubject={Boundary result for Hodge-Riemann positivity and absolute Hodge classes},
    pdfkeywords={Hodge conjecture, absolute Hodge classes, arithmetic positivity, Hodge-Riemann relations, Deligne absoluteness},
    colorlinks=true,
    linkcolor=black,
    urlcolor=blue,
    citecolor=black
}
\renewcommand*{\HyperDestNameFilter}[1]{en.#1}
\makeatletter
\let\auto@bib\@empty
\makeatother

\begin{document}

% ===================================================================
% TITLE PAGE
% ===================================================================

\title{Arithmetic Positivity and Absolute Hodge Classes:\\Why Hodge-Riemann Positivity Does Not\\Strengthen Deligne's Absoluteness}

\author{Lukas Geiger}
\affiliation{Independent Researcher, Bernau, Germany}

\date{\today}

\begin{abstract}
We introduce the concept of \textit{arithmetic positivity} for the Hodge Conjecture: a rational $(p,p)$-class $\alpha \in H^{2p}(X, \mathbb{Q})$ on a smooth projective variety $X/\mathbb{C}$ is called arithmetically positive if it simultaneously satisfies Hodge-Riemann positivity under all polarizations, Deligne's absoluteness under all automorphisms of $\mathbb{C}$, and Lefschetz compatibility. We prove that algebraic classes are arithmetically positive and record limited functoriality tests: finite pullback is controlled on the pullback ample subcone and becomes a full AP statement only under an explicit ample-cone comparison hypothesis; K\"unneth products are treated for primitive classes and product-generated ample cones; finite pushforward remains conjectural.

The central observation is a \textbf{boundary result}: the arithmetically positive cone coincides exactly with Deligne's absolute Hodge classes ($\mathrm{AP}^p(X) = \mathrm{AbsHodge}^p(X)$), because the Hodge-Riemann bilinear relations automatically guarantee the positivity conditions for any class of type $(p,p)$. This is a well-known consequence of the Hodge-Riemann relations; the contribution of this paper is to make the saturation boundary explicit and to delineate what Hodge-Riemann-based positivity methods cannot achieve. The natural attempt to strengthen Deligne's absoluteness through Hodge-Riemann positivity therefore produces no new constraints, and the Arithmetic Positivity Conjecture reduces to Deligne's question (1982): ``Are all absolute Hodge classes algebraic?'' Furthermore, the obstructions (O1) and (O3) in the No-Go formulation are vacuous for rational $(p,p)$-classes; only (O2), failure of absoluteness, is effective. We introduce the Galois-Hodge-Riemann spectrum as a visualization of the absoluteness test, illustrate the theoretical predictions numerically across selected test configurations, and discuss controlled test-family directions for finding genuinely stronger conditions beyond the Hodge-Riemann relations.

This paper is an expository contribution documenting a structural limitation; it does not contain new results beyond the observation that $\mathrm{AP} = \mathrm{AbsHodge}$.
\end{abstract}

\keywords{Hodge Conjecture, algebraic cycles, arithmetic positivity, absolute Hodge classes, Hodge-Riemann bilinear form, No-Go theorem}

\maketitle

\begin{center}
\footnotesize
ORCID: \url{https://orcid.org/0009-0005-7296-1534}\\
AI tools (Claude by Anthropic) were used for mathematical formalization, literature verification, and editorial refinement. All mathematical content was developed, verified, and approved by the author.
\end{center}


% ===================================================================
% 1. INTRODUCTION
% ===================================================================
\section{Introduction}
\label{sec:intro}

\subsection{The Hodge Conjecture}

Let $X$ be a smooth projective variety of dimension $n$ over $\mathbb{C}$. The Hodge decomposition
\begin{equation}
H^k(X, \mathbb{C}) = \bigoplus_{p+q=k} H^{p,q}(X)
\label{eq:hodge_decomp}
\end{equation}
expresses the $k$-th singular cohomology as a direct sum of Dolbeault cohomology groups. A class $\alpha \in H^{2p}(X, \mathbb{Q})$ is called a \textit{rational Hodge class} if its image in $H^{2p}(X, \mathbb{C})$ lies entirely in $H^{p,p}(X)$. Every algebraic cycle $Z$ of codimension $p$ defines a rational Hodge class $\mathrm{cl}(Z) \in H^{2p}(X, \mathbb{Q}) \cap H^{p,p}(X)$ via the cycle class map.

\begin{conjecture}[Hodge, 1950 \cite{Hodge1950}]
\label{conj:hodge}
Every rational Hodge class on a smooth projective variety over $\mathbb{C}$ is a $\mathbb{Q}$-linear combination of classes of algebraic cycles.
\end{conjecture}

Despite over seven decades of effort, the conjecture remains open in general. It is known to hold for:
\begin{itemize}
\item Divisors ($p = 1$): by the Lefschetz $(1,1)$-theorem.
\item Abelian varieties of dimension $\leq 5$: by work of Hazama~\cite{Hazama1994} and Moonen~\cite{Moonen1999}.
\item Products of elliptic curves and certain Fermat hypersurfaces.
\item Uniruled varieties (for degree-$2$ classes).
\end{itemize}

\subsection{Why construction fails}

All major approaches to the Hodge Conjecture share a common strategy: \textit{construct} the algebraic cycle representing a given Hodge class. This paper argues that this strategy is fundamentally misguided, and proposes an alternative: prove that non-algebraic Hodge classes are \textit{impossible} via a positivity obstruction.

The structural failures of existing approaches are:

\textbf{1. Direct algebraization.} Given a rational $(p,p)$-class, attempt to construct a subvariety whose fundamental class matches it. This fails because there is no canonical projection from Hodge cohomology to algebraic cycles, and the construction is not functorial under deformation.

\textbf{2. Variations of Hodge structure.} Use families $X_t$ and track how Hodge classes vary. This fails because Hodge classes can ``jump'' -- the condition of being type $(p,p)$ is not open in the Zariski topology -- and monodromy destroys naive stability arguments \cite{Voisin2002}.

\textbf{3. Motivic methods.} Replace varieties by motives and hope that Hodge classes are motivic. This fails because the category of motives is not fully constructed: Grothendieck's standard conjectures \cite{Grothendieck1969} remain open, and the relationship between Hodge classes and absolute Hodge classes is unresolved \cite{Andre1996}.

\textbf{4. Reduction mod $p$.} Reduce to characteristic $p$, invoke the Tate conjecture, and lift back. This fails because the bridge Hodge $\Rightarrow$ Tate $\Rightarrow$ Hodge is not closed: Frobenius invariance does not imply Hodge type, and lifting algebraic cycles is not unique \cite{Charles2013}.

As we show below, the positivity alternative proposed in this paper also encounters a fundamental limitation: the natural positivity conditions reduce to Deligne's absoluteness. Nevertheless, the analysis delineates a precise boundary for this class of approaches.

\subsection{The positivity alternative}

We observe that major breakthroughs in adjacent conjectures follow a different pattern:

\begin{center}
\resizebox{\columnwidth}{!}{%
\begin{tabular}{ll}
\toprule
Conjecture & Key positivity property \\
\midrule
Weil conjectures & Non-negativity of Frobenius eigenvalues \\
BSD (rank 1) & Non-negativity of N\'eron--Tate height \\
Tate (abelian var.) & Positivity via Faltings' isogeny theorem \\
\bottomrule
\end{tabular}}
\end{center}

In each case, a positivity property of the relevant arithmetic object plays a structural role -- though the proof techniques differ substantially (Deligne's proof of the Weil conjectures uses $\ell$-adic cohomology and the Lefschetz fixed-point theorem; Gross--Zagier explicitly constructs a Heegner point and computes its height). The Hodge Conjecture lacks: \textit{a positivity structure that forces algebraicity without constructing cycles}.

This paper introduces such a structure -- \textit{arithmetic positivity} -- combining Hodge-Riemann positivity under all polarizations with Deligne's absoluteness under all automorphisms of $\mathbb{C}$. The central finding, however, is a boundary/no-go observation: the Hodge-Riemann bilinear relations automatically guarantee the positivity conditions for any class of type $(p,p)$, so the arithmetically positive cone coincides with Deligne's absolute Hodge classes: $\mathrm{AP}^p(X) = \mathrm{AbsHodge}^p(X)$. The automatic saturation of Hodge-Riemann positivity for $(p,p)$-classes is well known to specialists (cf.\ Voisin~\cite{Voisin2002}, Chapter~6; Griffiths--Harris~\cite{GH1978}, Chapter~0--1); the purpose of this paper is to make the saturation boundary explicit and systematic within a unified positivity framework, and to clarify that no Hodge-Riemann-based strengthening of Deligne's absoluteness is possible. The Arithmetic Positivity Conjecture reduces to Deligne's question: ``Are all absolute Hodge classes algebraic?''

We regard this boundary result as informative: it delineates a precise limit for positivity-based approaches to the Hodge Conjecture and identifies directions for genuinely stronger conditions beyond Hodge-Riemann.

\subsection{Outline}

Section~\ref{sec:background} reviews the Hodge-Riemann bilinear form and absolute Hodge classes. Section~\ref{sec:arithmetic_positivity} introduces arithmetic positivity. Section~\ref{sec:easy_direction} proves that algebraic classes are arithmetically positive. Section~\ref{sec:functoriality} records limited functoriality tests and their hypotheses. Section~\ref{sec:nogo} formulates the No-Go theorem. Section~\ref{sec:abelian} proves the equivalence for abelian varieties. Section~\ref{sec:obstruction} identifies the precise obstruction. Section~\ref{sec:discussion} discusses the remaining gap and the relationship to Deligne's program.

This paper is self-contained; all main results depend only on definitions and proofs within this document. Connections to the broader FST programme are cited for context but are not logically required.


% ===================================================================
% 2. BACKGROUND
% ===================================================================
\section{Background}
\label{sec:background}

\subsection{Hodge-Riemann bilinear form}

Let $X$ be a smooth projective variety of dimension $n$, equipped with a K\"ahler form $\omega$ defining an ample class $L = [\omega] \in H^{1,1}(X) \cap H^2(X, \mathbb{Z})$. The \textit{Lefschetz operator} $\Lambda_L: H^k(X) \to H^{k+2}(X)$ is defined by $\Lambda_L(\alpha) = L \cup \alpha$. We use the convention $\Lambda_L(\alpha) = L \cup \alpha$; some authors reserve $\Lambda$ for the dual Lefschetz operator, while denoting the cup-product operation simply by $L$ or $L \wedge$ (Voisin~\cite{Voisin2002}; Griffiths--Harris~\cite{GH1978}).

\begin{theorem}[Hard Lefschetz]
\label{thm:hard_lefschetz}
For $k \leq n$, the map $\Lambda_L^{n-k}: H^k(X, \mathbb{Q}) \to H^{2n-k}(X, \mathbb{Q})$ is an isomorphism.
\end{theorem}

A class $\alpha \in H^k(X, \mathbb{Q})$ with $k \leq n$ is \textit{primitive} (with respect to $L$) if $\Lambda_L^{n-k+1}(\alpha) = 0$. The \textit{Hodge-Riemann bilinear form} on the primitive cohomology $P^k(X, \mathbb{Q}) \subset H^k(X, \mathbb{Q})$ is:
\begin{equation}
Q_L(\alpha, \beta) = \int_X \alpha \cup \beta \cup L^{n-k}\,.
\label{eq:HR_form}
\end{equation}

\begin{theorem}[Hodge-Riemann bilinear relations]
\label{thm:HR}
On the primitive subspace $P^{p,q}(X) = P^{p+q}(X) \cap H^{p,q}(X)$ with $p + q = k$, the Hermitian form
\begin{equation}
h_L(\alpha, \beta) = i^{p-q}\,(-1)^{k(k-1)/2}\, Q_L(\alpha, \bar{\beta})
\label{eq:HR_hermitian}
\end{equation}
is positive definite.
\end{theorem}

The Hodge-Riemann relations are the fundamental positivity result in Hodge theory. They provide definiteness \textit{relative to a fixed polarization} $L$. A central theme of this paper is that algebraicity requires positivity not just for one $L$, but for \textit{all} admissible polarizations and operations.

\subsection{Absolute Hodge classes}

Deligne \cite{Deligne1982} introduced a refinement of Hodge classes. An automorphism $\sigma \in \mathrm{Aut}(\mathbb{C})$ sends a smooth projective variety $X/\mathbb{C}$ to the conjugate variety $X^\sigma$ (same equations, coefficients twisted by $\sigma$). The comparison isomorphism provides:
\begin{equation}
H^{2p}_{\mathrm{dR}}(X/\mathbb{C}) \xrightarrow{\sim} H^{2p}_{\mathrm{dR}}(X^\sigma/\mathbb{C})\,,
\end{equation}
mapping the de~Rham class of $\alpha$ on $X$ to a class $\alpha^\sigma$ on $X^\sigma$.

\begin{definition}[Deligne]
\label{def:absolute}
A rational Hodge class $\alpha \in H^{2p}(X, \mathbb{Q}) \cap H^{p,p}(X)$ is \textit{absolute} (or \textit{absolutely Hodge}) if for every $\sigma \in \mathrm{Aut}(\mathbb{C})$, the conjugate class $\alpha^\sigma$ is again a Hodge class on $X^\sigma$.
\end{definition}

\begin{theorem}[Deligne \cite{Deligne1982}]
\label{thm:deligne_absolute}
\begin{enumerate}
\item Every algebraic class is absolute.
\item On abelian varieties, every Hodge class is absolute.
\end{enumerate}
\end{theorem}

Deligne's result shows that absoluteness is a \textit{necessary} condition for algebraicity. The question is whether it is sufficient. Our concept of arithmetic positivity \textit{attempts} to strengthen absoluteness by adding Hodge-Riemann positivity under all conjugations; as shown in Section~\ref{sec:discussion}, this attempt does not succeed (the additional conditions are automatic), but the analysis is informative.


% ===================================================================
% 3. ARITHMETIC POSITIVITY
% ===================================================================
\section{Arithmetic Positivity}
\label{sec:arithmetic_positivity}

We now introduce the central concept.

\subsection{Definition}

Let $X$ be a smooth projective variety of dimension $n$ over $\mathbb{C}$, and let $\mathrm{Amp}(X)$ denote the ample cone of $X$.

\begin{definition}[Arithmetic Positivity]
\label{def:arith_pos}
A rational Hodge class $\alpha \in H^{2p}(X, \mathbb{Q}) \cap H^{p,p}(X)$ is \textit{arithmetically positive} if it satisfies the following three conditions simultaneously:

When using the direct primitive notation below, we work in the middle-or-lower range $2p\leq n$; classes in complementary codimension are understood via Poincar\'e duality.  Likewise, the Lefschetz-compatibility condition is only imposed for cup-powers that still have a directly defined primitive component in the middle-or-lower range.

\begin{enumerate}
\item[\textbf{(AP1)}] \textbf{Universal Hodge-Riemann positivity.} For every ample class $L \in \mathrm{Amp}(X)$, the primitive component $\alpha_0$ of $\alpha$ (with respect to $L$) satisfies the \textit{signed} Hodge-Riemann inequality:
\begin{equation}
(-1)^p\, Q_L(\alpha_0, \alpha_0) \geq 0\,,
\label{eq:AP1}
\end{equation}
where $Q_L$ is the Hodge-Riemann form~\eqref{eq:HR_form}. The sign $(-1)^p$ is essential: the Hodge-Riemann bilinear relations (Theorem~\ref{thm:HR}) give $(-1)^p Q_L > 0$ on primitive $(p,p)$-classes, not $Q_L > 0$.

\item[\textbf{(AP2)}] \textbf{Absolute stability.} For every $\sigma \in \mathrm{Aut}(\mathbb{C})$:
\begin{enumerate}
\item[(a)] $\alpha^\sigma$ is a Hodge class on $X^\sigma$ (i.e., $\alpha$ is absolute).
\item[(b)] The conjugated class $\alpha^\sigma$ satisfies (AP1) on $X^\sigma$ with respect to every ample class $L^\sigma \in \mathrm{Amp}(X^\sigma)$.
\end{enumerate}

\item[\textbf{(AP3)}] \textbf{Lefschetz compatibility.} For all $0 \leq j \leq \lfloor n/2\rfloor - p$ and all ample $L$, let $\gamma_j = \Lambda_L^j \alpha \in H^{2(p+j)}(X)$ and let $\gamma_{j,0}$ denote its primitive component with respect to $L$ in $H^{2(p+j)}(X)$. Then:
\begin{equation}
(-1)^{p+j}\, Q_{L}(\gamma_{j,0},\, \gamma_{j,0}) \geq 0\,.
\label{eq:AP3}
\end{equation}
\end{enumerate}

The set of arithmetically positive classes in $H^{2p}(X, \mathbb{Q}) \cap H^{p,p}(X)$ is denoted $\mathrm{AP}^p(X)$.
\end{definition}

The full AP condition is \emph{a priori} stronger than Deligne's absoluteness alone; we will show in Section~8 that the two notions in fact coincide.

\begin{remark}
\label{rem:redundancy}
The conditions (AP1), (AP2b), and (AP3) are \textit{not} independent of (AP2a). As shown in Proposition~\ref{prop:absolute_implies_ap} below, the Hodge-Riemann bilinear relations automatically guarantee (AP1), (AP2b), and (AP3) for any class of type $(p,p)$. The only substantive condition is therefore (AP2a): Deligne's absoluteness. In particular, $\mathrm{AP}^p(X) = \mathrm{AbsHodge}^p(X)$ (Remark~\ref{rem:ap_equals_absolute}). The definition retains all three conditions for expository clarity: they make explicit the \textit{type} of positivity that algebraic classes enjoy, even though the Hodge-Riemann relations render them automatic for $(p,p)$-classes. Finding genuinely stronger conditions -- beyond the Hodge-Riemann relations -- that \textit{do} constrain beyond absoluteness remains an open problem (see Section~\ref{sec:discussion}).
\end{remark}

\subsection{The Arithmetic Positivity Conjecture}

\begin{conjecture}[Arithmetic Positivity Conjecture]
\label{conj:AP}
Let $X$ be a smooth projective variety over $\mathbb{C}$. Then:
\begin{equation}
\mathrm{AP}^p(X) = \mathrm{cl}\!\left(\mathrm{CH}^p(X)_\mathbb{Q}\right)\,,
\end{equation}
where $\mathrm{CH}^p(X)_\mathbb{Q}$ is the Chow group of codimension-$p$ cycles with rational coefficients, and $\mathrm{cl}$ is the cycle class map.
\end{conjecture}

\begin{proposition}
\label{prop:equivalence}
The following relations hold between the APC and the HC:
\begin{enumerate}
\item The Hodge Conjecture implies the Arithmetic Positivity Conjecture.
\item The ``hard direction'' of the APC ($\mathrm{AP}^p(X) \subseteq \mathrm{cl}(\mathrm{CH}^p(X)_\mathbb{Q})$) together with the statement that all rational Hodge classes are arithmetically positive ($H^{2p}(X, \mathbb{Q}) \cap H^{p,p}(X) \subseteq \mathrm{AP}^p(X)$) implies the HC.
\item Unconditionally: $\mathrm{cl}(\mathrm{CH}^p(X)_\mathbb{Q}) \subseteq \mathrm{AP}^p(X)$ (Theorem~\ref{thm:easy_direction} below).
\end{enumerate}
\end{proposition}

\begin{proof}
(1) If the HC holds, then every Hodge class is algebraic, and algebraic classes satisfy (AP1)--(AP3) (Theorem~\ref{thm:easy_direction} below). Thus $\mathrm{AP}^p(X) \subseteq H^{2p}(X, \mathbb{Q}) \cap H^{p,p}(X) = \mathrm{cl}(\mathrm{CH}^p(X)_\mathbb{Q})$, and the reverse inclusion holds by Theorem~\ref{thm:easy_direction}, giving $\mathrm{AP}^p(X) = \mathrm{cl}(\mathrm{CH}^p(X)_\mathbb{Q})$.

(2) This requires two separate assumptions. \textit{First}, suppose that the hard direction of the APC holds: $\mathrm{AP}^p(X) \subseteq \mathrm{cl}(\mathrm{CH}^p(X)_\mathbb{Q})$, i.e., every arithmetically positive class is algebraic. \textit{Second}, suppose that every rational Hodge class is arithmetically positive: $H^{2p}(X, \mathbb{Q}) \cap H^{p,p}(X) \subseteq \mathrm{AP}^p(X)$. Under both assumptions, every rational Hodge class $\alpha$ lies in $\mathrm{AP}^p(X)$ (by the second assumption) and hence in $\mathrm{cl}(\mathrm{CH}^p(X)_\mathbb{Q})$ (by the first assumption), so $\alpha$ is algebraic. This yields the HC.

Note that neither assumption alone suffices: the APC ($\mathrm{AP}^p = \mathrm{alg}^p$) shows that AP classes are algebraic, but without knowing that all Hodge classes are AP, one cannot conclude that all Hodge classes are algebraic.

(3) See Theorem~\ref{thm:easy_direction}.
\end{proof}

\begin{remark}
The APC alone does not directly imply the HC. The APC asserts $\mathrm{AP}^p = \mathrm{alg}^p$, while the HC asserts $\mathrm{Hodge}^p = \mathrm{alg}^p$. Since $\mathrm{AP}^p \subseteq \mathrm{Hodge}^p$ (arithmetic positivity imposes additional conditions beyond being a Hodge class), the APC gives $\mathrm{alg}^p \subseteq \mathrm{Hodge}^p$ (trivial) but not $\mathrm{Hodge}^p \subseteq \mathrm{alg}^p$ without the additional knowledge that all Hodge classes are arithmetically positive. The non-trivial content thus splits into two parts: (i)~show that non-AP Hodge classes do not exist (i.e., every Hodge class is AP), and (ii)~show that AP implies algebraic (the hard direction of APC).
\end{remark}

The non-trivial content of the APC is the ``hard direction'': $\mathrm{AP}^p(X) \subseteq \mathrm{cl}(\mathrm{CH}^p(X)_\mathbb{Q})$. This says that arithmetic positivity \textit{forces} algebraicity.


% ===================================================================
% 4. THE EASY DIRECTION
% ===================================================================
\section{Algebraic Classes are Arithmetically Positive}
\label{sec:easy_direction}

\begin{theorem}[Easy Direction]
\label{thm:easy_direction}
Let $X$ be a smooth projective variety over $\mathbb{C}$. Every class in the image of the cycle class map is arithmetically positive:
\begin{equation}
\mathrm{cl}\!\left(\mathrm{CH}^p(X)_\mathbb{Q}\right) \subseteq \mathrm{AP}^p(X)\,.
\end{equation}
\end{theorem}

\begin{proof}
Let $\alpha = \mathrm{cl}(Z)$ for some algebraic cycle $Z$ of codimension $p$ with rational coefficients. We verify (AP1)--(AP3).

\textit{(AP1):} Let $\alpha_0$ be the $L$-primitive component of $\alpha = \mathrm{cl}(Z)$. Since $\alpha$ is algebraic, it lies in $H^{p,p}(X) \cap H^{2p}(X, \mathbb{Q})$, and its primitive component $\alpha_0$ lies in $P^{p,p}(X) \cap H^{2p}(X, \mathbb{Q})$. By the Hodge-Riemann bilinear relations (Theorem~\ref{thm:HR}), the Hermitian form $h_L$ is positive definite on $P^{p,p}(X)$. For a rational $(p,p)$-class, $\bar{\alpha}_0 = \alpha_0$ (complex conjugation acts as the identity on classes of type $(p,p)$ with rational coefficients), $i^{p-p} = 1$, and $(-1)^{k(k-1)/2} = (-1)^{p(2p-1)}$ with $k = 2p$. Thus $h_L(\alpha_0, \alpha_0) = (-1)^{p(2p-1)}\, Q_L(\alpha_0, \alpha_0)$. Since $(-1)^{p(2p-1)} = (-1)^{2p^2 - p} = (-1)^p$ (as $(-1)^{2p^2} = 1$), the Hodge-Riemann relations give:
\begin{equation}
(-1)^p\, Q_L(\alpha_0, \alpha_0) > 0
\end{equation}
if $\alpha_0 \neq 0$. For $\alpha_0 = 0$ (i.e., when $\alpha$ has no primitive component), the inequality holds trivially. This argument uses only the Hodge-Riemann relations on $P^{p,p}(X)$ and does not require $\alpha_0$ itself to be the class of an effective cycle.

\textit{(AP2):} Part (a) is Deligne's theorem (Theorem~\ref{thm:deligne_absolute}): algebraic classes are absolute. Part (b) follows because $\sigma$ sends algebraic cycles on $X$ to algebraic cycles on $X^\sigma$ (algebraicity is a property of polynomial equations, which transform covariantly under $\sigma$), and algebraic classes on $X^\sigma$ satisfy (AP1) by the argument above.

\textit{(AP3):} The Hard Lefschetz theorem is compatible with the cycle class map: if $\alpha = \mathrm{cl}(Z)$, then $\Lambda_L^j \alpha = \mathrm{cl}(Z \cdot L^j)$, where $Z \cdot L^j$ is the intersection with $j$ copies of a hyperplane section. For the admissible range in Definition~\ref{def:arith_pos}, this is again an algebraic cycle of codimension $p + j$ with a directly defined $L$-primitive component $\gamma_{j,0}\in P^{p+j,\,p+j}(X)$. Since $\gamma_{j,0}$ is the primitive part of an algebraic class, the Hodge-Riemann relations give $(-1)^{p+j} Q_L(\gamma_{j,0}, \gamma_{j,0}) \geq 0$, by the same argument as in (AP1).
\end{proof}


% ===================================================================
% 5. FUNCTORIALITY
% ===================================================================
\section{Functoriality of Arithmetic Positivity}
\label{sec:functoriality}

For a full No-Go programme, one would want the arithmetically positive cone to be functorial under the standard operations of algebraic geometry. The results below record only the limited functoriality that is actually controlled here.

\begin{proposition}[Pullback -- finite morphisms]
\label{prop:pullback}
Let $f: Y \to X$ be a finite morphism of smooth projective varieties. If $\alpha \in \mathrm{AP}^p(X)$, then $f^*\alpha$ satisfies (AP1) for ample classes in the pullback subcone $f^*\mathrm{Amp}(X) \subset \mathrm{Amp}(Y)$, and satisfies (AP2)--(AP3) on that subcone. It is a full element of $\mathrm{AP}^p(Y)$ only under an explicit ample-cone comparison hypothesis, for example when the relevant positivity inequalities extend from $f^*\mathrm{Amp}(X)$ to all of $\mathrm{Amp}(Y)$ by density or by equality of the numerical ample cones.
\end{proposition}

\begin{proof}
Let $f: Y \to X$ be finite of degree $d$.

\textit{(AP1):} Since $f$ is finite, $f^*L_X$ is ample on $Y$ for every ample $L_X \in \mathrm{Amp}(X)$. For the primitive component $\alpha_0$ of $\alpha$ with respect to $L_X$, the projection formula gives:
\[
\int_Y f^*\alpha_0 \cup f^*\bar{\alpha}_0 \cup (f^*L_X)^{n-2p} = d \int_X \alpha_0 \cup \bar{\alpha}_0 \cup L_X^{n-2p},
\]
where $n = \dim X = \dim Y$. Since $d > 0$ and the right side satisfies $(-1)^p Q_{L_X}(\alpha_0, \alpha_0) \geq 0$ by assumption, we obtain $(-1)^p Q_{f^*L_X}(f^*\alpha_0, f^*\alpha_0) \geq 0$ on the pullback ample subcone. For a general ample class $L_Y \in \mathrm{Amp}(Y)$, no extension follows from finiteness alone: the primitive component of $f^*\alpha$ varies with $L_Y$, and the pullback subcone need not be dense in the ample cone of $Y$. If an additional ample-cone comparison hypothesis gives density or equality of the relevant cone, then continuity of the Hodge-Riemann form and of the Lefschetz decomposition extends the inequality to all ample classes.

\textit{(AP2):} Pullback commutes with $\sigma$-conjugation: $(f^*\alpha)^\sigma = (f^\sigma)^*(\alpha^\sigma)$, and $f^\sigma: Y^\sigma \to X^\sigma$ is again finite of the same degree. Since $\alpha^\sigma \in \mathrm{AP}^p(X^\sigma)$ by assumption, the result follows from the (AP1) argument applied to $f^\sigma$ and $\alpha^\sigma$.

\textit{(AP3):} The Lefschetz operator satisfies $\Lambda_{f^*L}^j \circ f^* = f^* \circ \Lambda_L^j$ (compatibility of pullback with cup product). Therefore $\Lambda_{f^*L}^j(f^*\alpha) = f^*(\Lambda_L^j \alpha)$, and the positivity of the primitive component follows from the same degree-$d$ argument applied to $\Lambda_L^j\alpha$.
\end{proof}

\begin{remark}
The degree formula $Q_{f^*L_X}(f^*\alpha_0, f^*\alpha_0) = \deg(f) \cdot Q_{L_X}(\alpha_0, \alpha_0)$ used in the proof relies on $f$ being a finite morphism (equivalently, generically finite with $\dim Y = \dim X$). For a generically finite morphism $f: Y \to X$ with $\dim Y = \dim X$ that is not finite, the projection formula still yields the degree relation on a dense open subset, but care is needed regarding the ramification locus; the statement of Proposition~\ref{prop:pullback} is therefore restricted to finite morphisms.
\end{remark}

\begin{remark}
For morphisms $f: Y \to X$ with $\dim Y > \dim X$, the pullback $f^*L_X$ is not ample on $Y$, and the primitive decomposition on $Y$ requires a relative Hodge-Riemann analysis. The functoriality of $\mathrm{AP}^p$ under general morphisms remains open.
\end{remark}

\begin{remark}[Density caveat]\label{rem:density-caveat}
The extension from $f^*\mathrm{Amp}(X)$ to all of $\mathrm{Amp}(Y)$ requires an explicit ample-cone comparison. A useful sufficient condition is that the pullback subcone has dense image in the relevant numerical ample cone, or that the numerical ample cones agree after pullback. This may hold in Picard-number-one situations, but it fails in general. Without such a hypothesis, Proposition~\ref{prop:pullback} is only a subcone functoriality statement.
\end{remark}

\begin{conjecture}[Pushforward -- finite morphisms]
\label{prop:pushforward}
Let $f: Y \to X$ be a finite morphism of smooth projective varieties (so $\dim Y = \dim X$). If $\alpha \in \mathrm{AP}^{p}(Y)$, then $f_*\alpha \in \mathrm{AP}^p(X)$.
\end{conjecture}

\begin{proof}[Evidence]
The projection formula $f_*f^* = \deg(f) \cdot \mathrm{id}$ on $H^*(X, \mathbb{Q})$ implies that $f_*$ is injective on the image of $f^*$ (up to scalar). For $\alpha \in \mathrm{AP}^p(Y)$ and any ample $L_X \in \mathrm{Amp}(X)$, the class $f^*L_X$ is ample on $Y$ (since $f$ is finite). The Hodge-Riemann form on $X$ is related to that on $Y$ via the adjunction $\int_X f_*\beta \cup \gamma = \int_Y \beta \cup f^*\gamma$. Combined with the degree formula, the sign conditions for $f_*\alpha$ follow from those of $\alpha$. The absolute stability transfers because $f^\sigma$ is again finite of the same degree. A complete proof requires verifying that the primitive decomposition on $X$ is compatible with $f_*$ applied to primitive classes on $Y$; this is not straightforward, since $f_*$ does not preserve primitivity in general.
\end{proof}

\begin{remark}
For morphisms $f: Y \to X$ with $\dim Y > \dim X$, the pushforward $f_*$ involves a Gysin-type argument relating the Hodge-Riemann form on $Y$ to that on $X$ via the Grothendieck-Riemann-Roch formula and the relative primitive decomposition. This analysis is not carried out here; the functoriality of $\mathrm{AP}^p$ under general pushforward remains open.
\end{remark}

\begin{proposition}[K\"unneth product -- primitive classes, product polarizations]
\label{prop:kunneth}
Let $\alpha \in \mathrm{AP}^p(X)$ and $\beta \in \mathrm{AP}^q(Y)$ be primitive with respect to some ample classes $L_X$ and $L_Y$ respectively. Then $\alpha \boxtimes \beta$ satisfies (AP1) for all product polarizations $L = L_X \boxtimes 1 + 1 \boxtimes L_Y$, and satisfies (AP2)--(AP3). In particular, $\alpha \boxtimes \beta \in \mathrm{AP}^{p+q}(X \times Y)$ when product polarizations generate $\mathrm{Amp}(X \times Y)$ (e.g., when $X$ or $Y$ has Picard number~$1$); for general product varieties, the extension to all ample classes remains open.
\end{proposition}

\begin{remark}
The extension to non-primitive classes requires a careful analysis of cross-terms in the Lefschetz decomposition on the product, which we do not carry out here. The result for primitive classes covers the structurally essential case.
\end{remark}

We first establish a key lemma on primitivity of exterior products.

\begin{lemma}[Primitivity of exterior products]
\label{lem:primitive_product}
Let $X$ and $Y$ be smooth projective varieties of dimensions $n$ and $m$ respectively, with ample classes $L_X \in \mathrm{Amp}(X)$ and $L_Y \in \mathrm{Amp}(Y)$. Set $L = L_X \boxtimes 1 + 1 \boxtimes L_Y \in \mathrm{Amp}(X \times Y)$. If $\alpha \in P^{2p}(X, \mathbb{Q})$ is primitive with respect to $L_X$ and $\beta \in P^{2q}(Y, \mathbb{Q})$ is primitive with respect to $L_Y$, then $\alpha \boxtimes \beta$ is primitive with respect to $L$ on $X \times Y$.
\end{lemma}

\begin{proof}
Set $k = p + q$ and $N = n + m = \dim(X \times Y)$. We must show $L^{N-2k+1} \cup (\alpha \boxtimes \beta) = 0$. Since $L = L_X \boxtimes 1 + 1 \boxtimes L_Y$, the binomial expansion gives for any $j \geq 0$:
\[
L^j \cup (\alpha \boxtimes \beta) = \sum_{a+b=j} \binom{j}{a} (L_X^a \cup \alpha) \boxtimes (L_Y^b \cup \beta).
\]
Each summand vanishes unless both factors are non-zero. Since $\alpha$ is $L_X$-primitive, $L_X^{n-2p+1} \cup \alpha = 0$, so $L_X^a \cup \alpha = 0$ for $a \geq n - 2p + 1$. Similarly, $L_Y^b \cup \beta = 0$ for $b \geq m - 2q + 1$. A summand with $a + b = j$ is non-zero only if $a \leq n - 2p$ and $b \leq m - 2q$, which requires $j \leq (n - 2p) + (m - 2q) = N - 2k$. Therefore $L^j \cup (\alpha \boxtimes \beta) = 0$ for all $j \geq N - 2k + 1$, which is precisely the primitivity condition.
\end{proof}

\begin{lemma}[Sign compatibility for exterior products]
\label{lem:sign_product}
Under the hypotheses of Lemma~\ref{lem:primitive_product}, with $k = p + q$:
\begin{equation}
\begin{aligned}
(-1)^k\, Q_L(\alpha \boxtimes \beta,\, \alpha \boxtimes \beta)
&= \binom{N-2k}{n-2p}
   \bigl[(-1)^p Q_{L_X}(\alpha, \alpha)\bigr] \\
&\quad\cdot \bigl[(-1)^q Q_{L_Y}(\beta, \beta)\bigr],
\end{aligned}
\label{eq:sign_kunneth}
\end{equation}
where $N = n + m$ and $Q_L$, $Q_{L_X}$, $Q_{L_Y}$ are the Hodge-Riemann forms~\eqref{eq:HR_form} on $X \times Y$, $X$, and $Y$ respectively.
\end{lemma}

\begin{proof}
By Lemma~\ref{lem:primitive_product}, $\alpha \boxtimes \beta$ is $L$-primitive, so the definition~\eqref{eq:HR_form} gives directly:
\[
Q_L(\alpha \boxtimes \beta,\, \alpha \boxtimes \beta) = \int_{X \times Y} (\alpha \boxtimes \beta)^2 \cup L^{N-2k}.
\]
Expanding $L^{N-2k}$ via the binomial formula as in Lemma~\ref{lem:primitive_product}, the only non-vanishing term occurs at $a = n - 2p$, $b = m - 2q$ (forced by $a \leq n-2p$, $b \leq m-2q$, and $a + b = N - 2k = (n-2p) + (m-2q)$). By the K\"unneth formula for integration:
\[
\begin{aligned}
\int_{X \times Y}(\alpha \boxtimes \beta)^2 \cup L^{N-2k}
&= \binom{N-2k}{n-2p}
   \left(\int_X \alpha^2 \cup L_X^{n-2p}\right)\\
&\quad\cdot
   \left(\int_Y \beta^2 \cup L_Y^{m-2q}\right).
\end{aligned}
\]
By definition~\eqref{eq:HR_form}, the integrals on $X$ and $Y$ are precisely $Q_{L_X}(\alpha, \alpha)$ and $Q_{L_Y}(\beta, \beta)$. Therefore:
\[
\begin{aligned}
Q_L(\alpha \boxtimes \beta,\, \alpha \boxtimes \beta)
&= \binom{N-2k}{n-2p}\,
   Q_{L_X}(\alpha, \alpha)\\
&\quad\cdot Q_{L_Y}(\beta, \beta).
\end{aligned}
\]
Multiplying both sides by $(-1)^k = (-1)^{p+q} = (-1)^p \cdot (-1)^q$ yields equation~\eqref{eq:sign_kunneth}.
\end{proof}

\begin{proof}[Proof of Proposition~\ref{prop:kunneth}]
We verify (AP1)--(AP3) for $\alpha \boxtimes \beta$ with $\alpha \in \mathrm{AP}^p(X)$ primitive and $\beta \in \mathrm{AP}^q(Y)$ primitive.

\textit{(AP1):} Let $L = L_X \boxtimes 1 + 1 \boxtimes L_Y$ for ample classes $L_X \in \mathrm{Amp}(X)$, $L_Y \in \mathrm{Amp}(Y)$. By Lemma~\ref{lem:primitive_product}, $\alpha \boxtimes \beta$ is $L$-primitive, so it equals its own primitive component. By Lemma~\ref{lem:sign_product}:
\[
\begin{aligned}
&(-1)^{p+q} Q_L(\alpha \boxtimes \beta,\, \alpha \boxtimes \beta)\\
&\quad= \binom{N-2k}{n-2p}
   \bigl[(-1)^p Q_{L_X}(\alpha, \alpha)\bigr]\\
&\quad\cdot \bigl[(-1)^q Q_{L_Y}(\beta, \beta)\bigr].
\end{aligned}
\]
The binomial coefficient $\binom{N-2k}{n-2p} > 0$, and both bracketed factors are $\geq 0$ by (AP1) for $\alpha$ and $\beta$. Hence $(-1)^{p+q}\, Q_L(\alpha \boxtimes \beta, \alpha \boxtimes \beta) \geq 0$. Ample classes of the form $L_X \boxtimes 1 + 1 \boxtimes L_Y$ form a subcone of $\mathrm{Amp}(X \times Y)$. Since the set $\{L \in \overline{\mathrm{Amp}}(X \times Y) : (-1)^{p+q} Q_L(\gamma_0, \gamma_0) \geq 0\}$ is closed (as $Q_L$ depends continuously on $L$) and contains this subcone, the inequality extends to the closure. For product varieties where $\mathrm{Amp}(X \times Y)$ is generated by pullback classes (e.g., when $X$ or $Y$ has Picard number~$1$), this suffices. In general, the extension to all of $\mathrm{Amp}(X \times Y)$ requires verifying the inequality for polarisations with mixed terms, which remains open.

\textit{(AP2):} For any $\sigma \in \mathrm{Aut}(\mathbb{C})$, $(\alpha \boxtimes \beta)^\sigma = \alpha^\sigma \boxtimes \beta^\sigma$ and $(X \times Y)^\sigma = X^\sigma \times Y^\sigma$. Since $\alpha^\sigma \in \mathrm{AP}^p(X^\sigma)$ and $\beta^\sigma \in \mathrm{AP}^q(Y^\sigma)$ by assumption, the (AP1) argument applies to $\alpha^\sigma \boxtimes \beta^\sigma$ on $X^\sigma \times Y^\sigma$.

\textit{(AP3):} The Lefschetz operator on the product satisfies $\Lambda_L^j(\alpha \boxtimes \beta) = \sum_{a+b=j} \binom{j}{a} (\Lambda_{L_X}^a \alpha) \boxtimes (\Lambda_{L_Y}^b \beta)$. For primitive $\alpha$ and $\beta$, the non-trivial terms have $a \leq n-2p$ and $b \leq m-2q$, and the positivity of each such term follows from (AP3) for $\alpha$ and $\beta$ together with the sign computation of Lemma~\ref{lem:sign_product}.

For non-primitive classes, the primitive decomposition on $X \times Y$ with respect to $L$ introduces cross-terms between different Lefschetz levels. The complete analysis of these cross-terms remains open; the above proof covers the structurally essential case of primitive classes.
\end{proof}


\begin{remark}[Barrier Lemma: HR-positivity is automatically saturated]
\label{rem:barrier-lemma}
Before proceeding to the No-Go theorem, we highlight the structural reason why any positivity-based strengthening of the Hodge Conjecture must exit the Hodge--Riemann framework. For a smooth projective variety $X$ and any ample class $L$, the Hodge--Riemann bilinear relations guarantee:
\begin{enumerate}
\item $(-1)^p Q_L(\alpha, \bar{\alpha}) > 0$ for every non-zero primitive $(p,p)$-class $\alpha$;
\item $Q_L$ is non-degenerate on $H^{p,p}_{\mathrm{prim}}(X)$.
\end{enumerate}
These conditions are \emph{automatically satisfied} by all rational $(p,p)$-classes---algebraic or not. Consequently, any positivity condition that is implied by the HR relations cannot distinguish algebraic from non-algebraic classes. This is the Barrier: the ``obvious'' strengthening (impose positivity under all polarisations, all automorphisms, all Lefschetz decompositions) collapses to the condition of being an absolute Hodge class (Theorem~\ref{thm:nogo}). The only escape routes require \emph{non-bilinear} positivity conditions (motivic convolution, Galois averaging, prismatic constraints) or information from \emph{outside} the Hodge--Riemann structure (o-minimality, derived categories).
\end{remark}

% ===================================================================
% 6. THE NO-GO THEOREM
% ===================================================================
\section{The No-Go Theorem}
\label{sec:nogo}

We now formulate the central result: the Hodge Conjecture restated as a positivity obstruction.

\begin{theorem}[Positivity Obstruction -- Conditional]
\label{thm:nogo}
Assume the Arithmetic Positivity Conjecture~\ref{conj:AP}. Let $X$ be a smooth projective variety over $\mathbb{C}$, and let $\alpha \in H^{2p}(X, \mathbb{Q}) \cap H^{p,p}(X)$ be a rational Hodge class that is \textbf{not} a $\mathbb{Q}$-linear combination of algebraic cycles. Then there exists at least one of the following obstructions:

\begin{enumerate}
\item[\textbf{(O1)}] \textbf{Polarization instability:} There exists an ample class $L \in \mathrm{Amp}(X)$ such that the primitive component $\alpha_0$ of $\alpha$ satisfies $(-1)^p\,Q_L(\alpha_0, \alpha_0) < 0$.

\item[\textbf{(O2)}] \textbf{Galois instability:} There exists $\sigma \in \mathrm{Aut}(\mathbb{C})$ such that either $\alpha^\sigma$ is not a Hodge class on $X^\sigma$, or $\alpha^\sigma$ violates (AP1) on $X^\sigma$.

\item[\textbf{(O3)}] \textbf{Lefschetz instability:} There exist an admissible $j \geq 1$ in the range of (AP3) and $L \in \mathrm{Amp}(X)$ such that $(-1)^{p+j}\,Q_L(\gamma_{j,0}, \gamma_{j,0}) < 0$, where $\gamma_{j,0}$ is the $L$-primitive component of $\Lambda_L^j \alpha \in H^{2(p+j)}(X)$.
\end{enumerate}
\end{theorem}

\begin{proof}
By contraposition: if $\alpha$ satisfies (AP1)--(AP3), then $\alpha \in \mathrm{AP}^p(X) = \mathrm{cl}(\mathrm{CH}^p(X)_\mathbb{Q})$ by the APC, contradicting the assumption that $\alpha$ is not algebraic.
\end{proof}

\begin{remark}[Structure of the obstruction]
\label{rem:obstruction_structure}
The three obstructions must be interpreted in light of the Hodge-Riemann relations.

\textbf{(O1) and (O3) are vacuous for $(p,p)$-classes.} The Hodge-Riemann bilinear relations (Theorem~\ref{thm:HR}) guarantee $(-1)^p Q_L(\alpha_0, \alpha_0) \geq 0$ for \textit{any} rational $(p,p)$-class $\alpha$, not only for algebraic ones. The primitive component $\alpha_0$ lies in $P^{p,p}(X)$, and the Hodge-Riemann form is positive definite there (with the sign $(-1)^p$). Similarly, for admissible AP3-powers the primitive component $\gamma_{j,0}$ of $\Lambda_L^j \alpha$ lies in $P^{p+j,p+j}(X)$, so $(-1)^{p+j} Q_L(\gamma_{j,0}, \gamma_{j,0}) \geq 0$ automatically. Therefore, obstructions (O1) and (O3) \textit{can never occur} for rational Hodge classes in the stated AP range.

\textbf{(O2) is the only effective obstruction.} The only way a Hodge class can fail arithmetic positivity is through (O2a): failure of absoluteness ($\alpha^\sigma$ is not Hodge on $X^\sigma$ for some $\sigma$). Part (O2b) -- failure of Hodge-Riemann positivity on $X^\sigma$ -- is also vacuous, since $\alpha^\sigma \in H^{p,p}(X^\sigma)$ automatically satisfies the Hodge-Riemann inequalities.

Consequently, the No-Go theorem reduces to Deligne's observation: non-absolute Hodge classes (if they exist) cannot be algebraic.
\end{remark}

\subsection{The unconditional statement}

While Theorem~\ref{thm:nogo} is conditional on the APC, we can formulate an unconditional partial result:

\begin{proposition}[Unconditional Obstruction]
\label{prop:unconditional}
Let $\alpha \in H^{2p}(X, \mathbb{Q}) \cap H^{p,p}(X)$ be a rational Hodge class. If $\alpha$ fails to satisfy (AP2a) -- i.e., there exists $\sigma \in \mathrm{Aut}(\mathbb{C})$ such that $\alpha^\sigma \notin H^{p,p}(X^\sigma)$ -- then $\alpha$ is not algebraic.
\end{proposition}

\begin{proof}
This is Deligne's result (Theorem~\ref{thm:deligne_absolute}, part~1, contrapositive): algebraic classes are absolute.
\end{proof}

As shown in Section~\ref{sec:discussion} (Remark~\ref{rem:ap_equals_absolute}), the Hodge-Riemann positivity conditions (AP1), (AP2b), and (AP3) are automatically satisfied by any absolute Hodge class. The content of the APC beyond Proposition~\ref{prop:unconditional} is therefore identical to Deligne's question: \textit{are all absolute Hodge classes algebraic?} The ``gap'' between absolute Hodge classes and algebraic classes is not narrowed by the positivity conditions of this paper.


% ===================================================================
% 7. THE CASE OF ABELIAN VARIETIES
% ===================================================================
\section{Abelian Varieties}
\label{sec:abelian}

Abelian varieties provide the most favorable setting for the arithmetic positivity approach, because Deligne's theorem (all Hodge classes are absolute) eliminates obstruction (O2a).

\begin{proposition}[Equivalence for abelian varieties]
\label{thm:abelian}
For abelian varieties, the Arithmetic Positivity Conjecture~\ref{conj:AP} is equivalent to the Hodge Conjecture.
\end{proposition}

\begin{proof}
Let $A$ be an abelian variety of dimension $g$, and let $\alpha \in H^{2p}(A, \mathbb{Q}) \cap H^{p,p}(A)$.

By Deligne's theorem (Theorem~\ref{thm:deligne_absolute}, part~2), every Hodge class on an abelian variety is absolute. Proposition~\ref{prop:absolute_implies_ap} then gives
\[
H^{2p}(A,\mathbb{Q})\cap H^{p,p}(A)=\mathrm{AbsHodge}^p(A)=\mathrm{AP}^p(A).
\]
Therefore the APC for $A$,
\[
\mathrm{AP}^p(A)=\mathrm{cl}(\mathrm{CH}^p(A)_{\mathbb{Q}}),
\]
is literally the Hodge Conjecture for $A$,
\[
H^{2p}(A,\mathbb{Q})\cap H^{p,p}(A)
=\mathrm{cl}(\mathrm{CH}^p(A)_{\mathbb{Q}}).
\]
This proves the equivalence. It does not prove the Hodge Conjecture for general abelian varieties; it only shows that arithmetic positivity adds no extra condition once Deligne's absoluteness theorem is available.
\end{proof}

\begin{remark}[Known positive cases]
For abelian varieties where all Hodge classes are known to be algebraic (for example several low-dimensional, CM, or product cases in the literature), the equivalent APC statement is verified as well. In general, however, Deligne's theorem supplies absoluteness rather than algebraicity; the remaining problem is still to show that the absolute Hodge classes are represented by algebraic cycles.
\end{remark}


% ===================================================================
% 8. THE PRECISE OBSTRUCTION
% ===================================================================
\section{The Precise Obstruction}
\label{sec:obstruction}

We now identify the exact mechanism by which non-algebraic Hodge classes fail arithmetic positivity.

\subsection{The Galois-Hodge-Riemann form}

\begin{definition}[Galois-Hodge-Riemann form]
\label{def:galois_HR}
For $\sigma \in \mathrm{Aut}(\mathbb{C})$ and $L \in \mathrm{Amp}(X)$, define:
\begin{equation}
Q_L^\sigma(\alpha) := Q_{L^\sigma}(\alpha^\sigma, \alpha^\sigma)\,.
\label{eq:galois_HR}
\end{equation}
The \textit{Galois-Hodge-Riemann spectrum} of $\alpha$ is:
\begin{equation}
\begin{aligned}
\mathrm{Spec}_{\mathrm{GHR}}(\alpha)
:= \{\, &(-1)^p Q_L^\sigma(\alpha) :\\
&\sigma \in \mathrm{Aut}(\mathbb{C}),\,
  L \in \mathrm{Amp}(X)\,\} \subset \mathbb{R}.
\end{aligned}
\end{equation}
\end{definition}

\begin{proposition}[Algebraic classes have non-negative GHR spectrum]
\label{prop:ghr_positive}
If $\alpha = \mathrm{cl}(Z)$ for some algebraic cycle $Z$, then:
\begin{equation}
\mathrm{Spec}_{\mathrm{GHR}}(\alpha) \subseteq [0, \infty)\,.
\end{equation}
\end{proposition}

\begin{proof}
Algebraic cycles transform covariantly under $\sigma$: $\mathrm{cl}(Z)^\sigma = \mathrm{cl}(Z^\sigma)$, and $Z^\sigma$ is an algebraic cycle on $X^\sigma$. The primitive component of $\mathrm{cl}(Z^\sigma)$ lies in $P^{p,p}(X^\sigma)$, so the Hodge-Riemann relations give $(-1)^p\, Q_{L^\sigma}(\mathrm{cl}(Z^\sigma)_0, \mathrm{cl}(Z^\sigma)_0) \geq 0$, which is the required non-negativity of the GHR spectrum.
\end{proof}

\begin{conjecture}[GHR Obstruction Conjecture]
\label{conj:ghr}
If $\alpha$ is a rational Hodge class that is \textbf{not} algebraic, then:
\begin{equation}
\mathrm{Spec}_{\mathrm{GHR}}(\alpha) \cap (-\infty, 0) \neq \emptyset\,.
\end{equation}
That is, there exist $\sigma$ and $L$ such that $Q_L^\sigma(\alpha) < 0$.
\end{conjecture}

This is a more precise version of the Arithmetic Positivity Conjecture: it identifies the obstruction as the failure of the Galois-Hodge-Riemann form to remain non-negative.

\begin{remark}[Logical status of the GHR Obstruction Conjecture]
\label{rem:ghr_status}
The GHR Obstruction Conjecture must be interpreted carefully in light of $\mathrm{AP}^p = \mathrm{AbsHodge}^p$ (Remark~\ref{rem:ap_equals_absolute}). If $\alpha$ is an absolute Hodge class, then $(-1)^p Q_L^\sigma(\alpha) \geq 0$ for all $\sigma$ and $L$ by Proposition~\ref{prop:absolute_implies_ap}. Therefore, the conjecture can only ``fire'' for non-algebraic Hodge classes that are \textit{not absolute} -- in which case the obstruction is already witnessed by (AP2a) (failure of absoluteness), not by negativity of $Q_L^\sigma$. If non-algebraic absolute Hodge classes exist (i.e., if Deligne's question has a negative answer), then the GHR Obstruction Conjecture would be \textit{false} for such classes, since their GHR spectrum is automatically non-negative. The conjecture is therefore equivalent to the assertion that non-algebraic absolute Hodge classes do not exist -- which is precisely Deligne's question.
\end{remark}

\begin{remark}[Falsifiability of the GHR diagnostic]
\label{rem:ghr-falsifiability}
The GHR spectrum can be used as a diagnostic tool with clear falsification criteria:
\begin{enumerate}
\item \emph{Positive test:} If a class $\alpha$ has $(-1)^p Q_L^\sigma(\alpha^\sigma) < 0$ for some $\sigma \in \mathrm{Aut}(\mathbb{C})$ and some $L$, then $\alpha$ is provably non-algebraic (by Theorem~\ref{thm:easy_direction}).
\item \emph{Null test:} If $(-1)^p Q_L^\sigma(\alpha^\sigma) \geq 0$ for all $\sigma$ and all $L$, then $\alpha$ is arithmetically positive---hence absolute Hodge---but not necessarily algebraic. This is the regime where the GHR diagnostic is \emph{silent}.
\item \emph{Structural test:} The GHR spectrum detects non-algebraicity only for classes that are \emph{not} absolute Hodge. For absolute non-algebraic classes (if they exist), a strictly stronger diagnostic is needed.
\end{enumerate}
The GHR spectrum thus functions as a \emph{necessary condition filter}: it can certify non-algebraicity but cannot certify algebraicity. This asymmetry mirrors the Pattern~A architecture across the FST programme, where positivity certificates exclude alternatives but do not construct solutions.
\end{remark}

\begin{example}[GHR spectrum of the diagonal class]
\label{ex:diagonal}
Let $X = \mathbb{P}^n$ and consider the diagonal embedding $\Delta: \mathbb{P}^n \hookrightarrow \mathbb{P}^n \times \mathbb{P}^n$. The class $[\Delta] \in H^{2n}(\mathbb{P}^n \times \mathbb{P}^n, \mathbb{Q})$ is algebraic (it is the cycle class of the diagonal). By Theorem~\ref{thm:easy_direction}, $[\Delta]$ is arithmetically positive. Its K\"unneth decomposition is $[\Delta] = \sum_{k=0}^n h^k \boxtimes h^{n-k}$, where $h$ is the hyperplane class. Since $\mathbb{P}^n$ has no non-trivial automorphisms of $\mathbb{C}$ acting on its cohomology (the Hodge structure is pure Tate), every $\sigma \in \mathrm{Aut}(\mathbb{C})$ fixes $[\Delta]^\sigma = [\Delta]$. The self-intersection number is $\int_{X \times X} [\Delta]^2 = \chi(\mathbb{P}^n) = n + 1$ (the Euler characteristic). Since $[\Delta]$ is algebraic, its primitive component $[\Delta]_0$ satisfies $(-1)^n Q_L([\Delta]_0, [\Delta]_0) \geq 0$ by the Hodge-Riemann relations. We note that $[\Delta]$ is \emph{not} primitive: $L \cup [\Delta] \neq 0$ in general, so the primitive decomposition is non-trivial. This example is ``trivially positive'' because algebraic, but it illustrates the computation: for a hypothetical non-algebraic Hodge class on a more complex variety, one would seek $\sigma$ and $L$ where the GHR form turns negative.
\end{example}

\begin{example}[Non-trivial GHR spectrum: CM endomorphism class]
\label{ex:cm_endomorphism}
Consider the elliptic curve $E: y^2 = x^3 - x$ over $\mathbb{Q}(i)$, which has complex multiplication by $\mathbb{Z}[i]$, with $\mathrm{End}(E) \otimes \mathbb{Q} = \mathbb{Q}(i)$. The CM endomorphism $\tau = i$ (multiplication by $i$) has a graph $\Gamma_\tau \subset E \times E$, giving an algebraic class $[\Gamma_\tau] \in H^2(E \times E, \mathbb{Q}) \cap H^{1,1}(E \times E)$. This class lies \emph{outside} the Lefschetz ring: the Lefschetz ring of $E \times E$ (generated by the two fibre classes $[E \times \{0\}]$, $[\{0\} \times E]$ and the diagonal $[\Delta]$) does not contain $[\Gamma_\tau]$, since $\tau$ acts non-trivially on $H^{1,0}(E)$.

The GHR spectrum of $[\Gamma_\tau]$ is non-trivial. Let $\sigma \in \mathrm{Aut}(\mathbb{C})$ be complex conjugation restricted to $\mathbb{Q}(i)$, so $\sigma(i) = -i$. Then $E^\sigma \cong E$ (since $y^2 = x^3 - x$ has rational coefficients), but the CM endomorphism transforms as $\tau^\sigma = \sigma(i) = -i$, so $[\Gamma_\tau]^\sigma = [\Gamma_{-\tau}]$. Let $L = L_1 \boxtimes 1 + 1 \boxtimes L_1$ be the product polarisation, where $L_1$ is the principal polarisation on $E$. Then:
\begin{align}
(-1)^1 Q_L^\mathrm{id}([\Gamma_\tau])
&= (-1)^1 Q_L([\Gamma_i], [\Gamma_i]) \notag\\
&= 1 + |\tau|^2 = 2, \label{eq:cm_id}\\
(-1)^1 Q_L^\sigma([\Gamma_\tau])
&= (-1)^1 Q_{L^\sigma}([\Gamma_{-i}], [\Gamma_{-i}]) \notag\\
&= 1 + |-\tau|^2 = 2. \label{eq:cm_sigma}
\end{align}
Both values are positive, as guaranteed by Theorem~\ref{thm:easy_direction}. However, the \emph{intermediate} computation differs: the class $[\Gamma_i]$ and $[\Gamma_{-i}]$ have different K\"unneth decompositions in $H^1(E) \otimes H^1(E)$, reflecting the fact that $\sigma$ permutes the eigenspaces of the CM action on $H^1(E, \mathbb{C}) = H^{1,0}(E) \oplus H^{0,1}(E)$. Explicitly, $\tau^*$ acts on $H^{1,0}(E)$ by multiplication by $i$ and on $H^{0,1}(E)$ by $-i$; under $\sigma$, these eigenvalues are exchanged.

The essential point is that the GHR spectrum $\{(-1)^p Q_L^\sigma([\Gamma_\tau])\}_{\sigma, L}$ takes \emph{different values for different $\sigma$} (due to the varying K\"unneth structure), yet remains uniformly non-negative. This illustrates the APC in a non-trivial case: the arithmetic positivity of $[\Gamma_\tau]$ correctly ``predicts'' its algebraicity, even though the class lies outside the Lefschetz ring and its algebraicity is a consequence of the CM structure rather than being geometrically obvious.
\end{example}

\subsection{Why the obstruction is natural}

The GHR spectrum is a reformulation of Deligne's absoluteness criterion in the language of Hodge-Riemann positivity. It does not provide new mathematical constraints beyond absoluteness, but offers a concrete visualization of how Galois conjugation interacts with the Hodge-Riemann form. The reformulation is natural for the following reason. The Hodge-Riemann form $Q_L$ depends on both:
\begin{itemize}
\item the \textit{complex structure} of $X$ (through the Hodge decomposition), and
\item the \textit{algebraic structure} (through the polarization $L$).
\end{itemize}

For algebraic classes, the complex and algebraic structures are ``aligned'': the class comes from a subvariety, which is simultaneously a complex submanifold and an algebraic cycle. The Galois action $\sigma$ can ``misalign'' the complex and algebraic structures, but for algebraic classes, the cycle $Z^\sigma$ ``re-aligns'' them on $X^\sigma$.

For a \textit{non-algebraic} Hodge class, there is no cycle to perform this re-alignment. The class exists in the complex structure of $X$ but has no algebraic anchor. When $\sigma$ twists the complex structure, the positivity of the Hodge-Riemann form may fail, because there is nothing to compensate for the twist.

This is the precise mechanism of the obstruction: \textit{algebraic classes are rigid enough to survive all Galois twists with their positivity intact, while transcendental classes are not.} We emphasize, however, that by Remark~\ref{rem:ap_equals_absolute} and Remark~\ref{rem:obstruction_structure}, the Hodge-Riemann positivity plays no independent role here: the sole effective obstruction is failure of absoluteness (O2a). The ``alignment'' narrative describes the mechanism of absoluteness, not of Hodge-Riemann positivity.


% ===================================================================
% 9. DISCUSSION
% ===================================================================
\section{Discussion}
\label{sec:discussion}

\subsection{The remaining gap}

The paper establishes:
\begin{enumerate}
\item Algebraic classes $\Rightarrow$ arithmetically positive (Theorem~\ref{thm:easy_direction}, unconditional).
\item $\mathrm{AP}^p(X) = \mathrm{AbsHodge}^p(X)$ (Remark~\ref{rem:ap_equals_absolute}, unconditional).
\item Limited functoriality tests for finite pullback subcones and K\"unneth products of primitive classes under product-polarization hypotheses (Section~\ref{sec:functoriality}).
\item Identification of the Galois-Hodge-Riemann spectrum as a visualization of the absoluteness criterion (Section~\ref{sec:obstruction}).
\end{enumerate}

As a consequence of point~2, the remaining gap is precisely Deligne's question:

\begin{quote}
\textit{Are all absolute Hodge classes algebraic?}
\end{quote}

This is a long-standing open problem. The positivity framework does not resolve it, but provides a new angle: the No-Go formulation and the GHR spectrum make explicit which Galois operations could witness failure of absoluteness. They do not distinguish algebraic from non-algebraic absolute Hodge classes if such classes exist.

\subsection{Relation to Deligne's program}

Deligne's framework of absolute Hodge classes \cite{Deligne1982} provides condition (AP2a). The present paper initially intended to strengthen Deligne's absoluteness by adding:

\begin{enumerate}
\item \textbf{Hodge-Riemann positivity under conjugation} (AP2b): Not just the Hodge \textit{type} but the Hodge-Riemann \textit{sign} must be preserved.

\item \textbf{Universal polarization} (AP1 for all $L$): All ample classes must give non-negative Hodge-Riemann form.
\end{enumerate}

However, Proposition~\ref{prop:absolute_implies_ap} and Remark~\ref{rem:ap_equals_absolute} show that these conditions are \textit{automatically} satisfied by any absolute Hodge class, via the Hodge-Riemann bilinear relations. The conditions (AP1), (AP2b), and (AP3) therefore do not restrict beyond Deligne's absoluteness. The question ``Are all absolute Hodge classes algebraic?'' remains precisely Deligne's motivating problem, and the APC is equivalent to it.

\begin{proposition}
\label{prop:absolute_implies_ap}
On a smooth projective variety $X/\mathbb{C}$, every absolute Hodge class is arithmetically positive:
\begin{equation}
\mathrm{AbsHodge}^p(X) \subseteq \mathrm{AP}^p(X)\,.
\end{equation}
\end{proposition}

\begin{proof}
Let $\alpha \in H^{2p}(X, \mathbb{Q}) \cap H^{p,p}(X)$ be an absolute Hodge class.

\textit{(AP1):} The primitive component $\alpha_0$ of $\alpha$ with respect to any ample $L$ lies in $P^{p,p}(X) \cap H^{2p}(X, \mathbb{Q})$. For rational $(p,p)$-classes, $\bar{\alpha}_0 = \alpha_0$ and $i^{p-p} = 1$, so $h_L(\alpha_0, \alpha_0) = (-1)^{p(2p-1)} Q_L(\alpha_0, \alpha_0) = (-1)^p Q_L(\alpha_0, \alpha_0)$. The Hodge-Riemann bilinear relations (Theorem~\ref{thm:HR}) give $h_L > 0$, hence $(-1)^p Q_L(\alpha_0, \alpha_0) > 0$ for $\alpha_0 \neq 0$. For $\alpha_0 = 0$, the inequality holds trivially.

\textit{(AP2):} Part (a) holds by assumption ($\alpha$ is absolute). For part (b), since $\alpha$ is absolute, $\alpha^\sigma$ is a Hodge class on $X^\sigma$ for every $\sigma \in \mathrm{Aut}(\mathbb{C})$. In particular, $\alpha^\sigma \in H^{p,p}(X^\sigma) \cap H^{2p}(X^\sigma, \mathbb{Q})$, and its $L^\sigma$-primitive component lies in $P^{p,p}(X^\sigma)$. The Hodge-Riemann relations on $X^\sigma$ then give $(-1)^p Q_{L^\sigma}(\alpha_0^\sigma, \alpha_0^\sigma) \geq 0$, exactly as in (AP1).

\textit{(AP3):} For admissible $j$ in the range of Definition~\ref{def:arith_pos}, the class $\Lambda_L^j \alpha \in H^{2(p+j)}(X)$ is a Hodge class of type $(p+j, p+j)$, and its primitive component lies in $P^{p+j, p+j}(X)$. The Hodge-Riemann relations give $(-1)^{p+j} Q_L(\gamma_{j,0}, \gamma_{j,0}) \geq 0$.
\end{proof}

\begin{remark}[Critical observation]
\label{rem:ap_equals_absolute}
Proposition~\ref{prop:absolute_implies_ap} shows that $\mathrm{AbsHodge}^p(X) \subseteq \mathrm{AP}^p(X)$. The reverse inclusion $\mathrm{AP}^p(X) \subseteq \mathrm{AbsHodge}^p(X)$ holds by definition, since condition (AP2a) requires absoluteness. Therefore:
\begin{equation}
\mathrm{AP}^p(X) = \mathrm{AbsHodge}^p(X)\,.
\end{equation}
In other words, the conditions (AP1), (AP2b), and (AP3) do \textit{not} impose any additional restriction beyond Deligne's absoluteness. The Hodge-Riemann bilinear relations \textit{automatically} guarantee the positivity conditions for any class of type $(p,p)$.

This has a profound consequence for the APC: the Arithmetic Positivity Conjecture $\mathrm{AP}^p(X) = \mathrm{cl}(\mathrm{CH}^p(X)_\mathbb{Q})$ is \textbf{equivalent} to the statement that every absolute Hodge class is algebraic -- which is precisely Deligne's original question \cite{Deligne1982}. The positivity framework introduced in this paper therefore does not produce a genuinely new conjecture, but rather provides a new \textit{perspective} on Deligne's question through the language of Hodge-Riemann positivity and No-Go obstructions.

The GHR obstruction (Definition~\ref{def:galois_HR}) remains meaningful as a concrete visualization of the absoluteness criterion: it makes explicit \textit{which} Galois conjugation would have to fail for a non-algebraic absolute Hodge class. However, the core mathematical content reduces to: \textit{are all absolute Hodge classes algebraic?}
\end{remark}

Combining Proposition~\ref{prop:absolute_implies_ap} with Theorem~\ref{thm:easy_direction}, the unconditional chain is:
\begin{equation}
\begin{aligned}
\text{Algebraic}
&\xRightarrow{\text{Thm.~\ref{thm:easy_direction}}}
  \text{AP} = \text{Absolute}\\
&\xhookrightarrow{\text{definition}} \text{Hodge}\,,
\end{aligned}
\end{equation}
where the equality $\text{AP} = \text{Absolute}$ is Remark~\ref{rem:ap_equals_absolute}.  The HC requires the composite reverse inclusion $\text{Hodge}\subseteq\text{Algebraic}$, which splits into two logically distinct questions: whether all Hodge classes are absolute, and whether all absolute Hodge classes are algebraic.  Deligne's theorem supplies the first reverse inclusion for abelian varieties, but the second reverse inclusion is precisely the remaining algebraicity problem.

\subsection{Connection to the broader program}

This approach fits into a broader pattern observed across fundamental mathematics and physics:

The approach fits into a family of results where positivity properties play a structural role:
\begin{itemize}
\item \textit{Weil conjectures:} The non-negativity of Frobenius eigenvalues is the \emph{statement}; the proof technique (Deligne~\cite{Deligne1974}) uses $\ell$-adic cohomology and the Lefschetz fixed-point theorem.
\item \textit{BSD (rank 1):} The non-negativity of the N\'eron--Tate height (Gross--Zagier~\cite{GrossZagier1986}) is a \emph{consequence} of the explicit Heegner-point construction.
\item \textit{Hodge:} The GHR spectrum reformulates Deligne's absoluteness in the language of Hodge-Riemann positivity (this paper); the core mathematical content reduces to Deligne's question.
\end{itemize}
We emphasize that the analogy is structural, not methodological: the proof techniques in these three settings are fundamentally different. Previous work by the author \cite{GeigerRH2026c, FST-Unified2026} explored related positivity-based approaches in other settings.

\begin{remark}[Saturation phenomenon]
\label{rem:meta-pattern-saturation}
The structural lesson of this paper -- that the ``obvious'' positivity strengthening is automatically saturated by the Hodge-Riemann relations -- is not specific to the Hodge Conjecture. A similar phenomenon occurs for height positivity in the BSD context (saturated by Gross--Zagier + Kolyvagin for rank $\leq 1$; the ``next'' positivity for rank $\geq 2$ requires objects that do not yet exist). In both cases, progress requires \emph{exiting} the saturated framework by introducing non-bilinear or arithmetic conditions.
\end{remark}

\subsection{Limitations}

\begin{enumerate}
\item \textbf{The hard direction is open.} This paper does not prove the Hodge Conjecture. It reformulates it as a positivity statement and identifies the precise obstruction, but the proof that arithmetic positivity forces algebraicity remains the central challenge. Moreover, the APC ($\mathrm{AP}^p = \mathrm{alg}^p$) does not by itself imply the HC ($\mathrm{Hodge}^p = \mathrm{alg}^p$): one additionally needs that all Hodge classes are arithmetically positive (see Proposition~\ref{prop:equivalence}).

\item \textbf{Partial sketches remain in Section~\ref{sec:functoriality}.} The K\"unneth functoriality is fully proved for primitive classes (Proposition~\ref{prop:kunneth}, Lemmas~\ref{lem:primitive_product}--\ref{lem:sign_product}), but the extension to non-primitive classes and the pushforward for fibred morphisms remain sketches requiring careful treatment of cross-terms in the primitive decomposition.

\item \textbf{The AP cone equals the absolute Hodge cone (the central limitation).} As shown in Remark~\ref{rem:ap_equals_absolute}, $\mathrm{AP}^p(X) = \mathrm{AbsHodge}^p(X)$, because the Hodge-Riemann bilinear relations automatically guarantee all positivity conditions for $(p,p)$-classes. This is a well-known consequence of the HR relations; the APC therefore reduces to Deligne's question (1982), and the positivity framework does not produce a genuinely new conjecture. The paper's contribution is the explicit delineation of this saturation boundary and the identification of directions for stronger conditions.

\item \textbf{Remaining abelian-variety gap.} For abelian varieties, Deligne's theorem identifies all Hodge classes with absolute Hodge classes, so AP coincides with the full Hodge-class space. The unsolved part is not absoluteness but algebraicity of those absolute classes; this is known only in special families and remains open in general.
\end{enumerate}

\subsection{Outlook}

Four directions are immediate:

\begin{enumerate}
\item \textbf{Verify the GHR Obstruction Conjecture for known non-algebraic classes.} A test case is Voisin's counterexample to the Hodge conjecture for K\"ahler manifolds \cite{Voisin2002b}. However, as noted in Remark~\ref{rem:ghr_status}, the conjecture can only detect non-absoluteness; for absolute Hodge classes that might not be algebraic, the GHR spectrum is automatically non-negative.

\item \textbf{Strengthen the AP conditions.} Since $\mathrm{AP}^p = \mathrm{AbsHodge}^p$ (Remark~\ref{rem:ap_equals_absolute}), the current definition of arithmetic positivity does not go beyond Deligne's absoluteness. One may seek \textit{additional} positivity conditions -- beyond those implied by the Hodge-Riemann relations -- that could narrow the gap between absolute Hodge classes and algebraic classes. Candidates include conditions involving higher-order Hodge-Riemann forms, $p$-adic Hodge theory, or motivic constraints.

\item \textbf{Develop the category of arithmetically positive classes} as a substitute for (part of) the theory of motives, with functorial operations and compatibility with the standard conjectures.

\item \textbf{The diagnostic quadratic locus.} The Hodge-Riemann part of the AP definition is expressed by \emph{quadratic} inequalities (the Hodge-Riemann form is bilinear, so the positivity conditions are quadratic in $\alpha$). Since the full AP locus is $\mathrm{AbsHodge}^p(X)$ by Remark~\ref{rem:ap_equals_absolute}, this quadratic locus is not a separate algebraicity criterion. It is a diagnostic shadow of the positivity conditions: it is closed under positive scaling and symmetric, but its boundary records only where the Hodge-Riemann visualization becomes degenerate. Understanding that boundary may clarify the limitations of HR-based diagnostics, not solve the hard direction of the APC.
\end{enumerate}


\begin{remark}[Arakelov bridge to BSD]
\label{rem:arakelov-bridge}
The Hodge--Riemann form $(-1)^p Q_L$ can be interpreted as positivity at the archimedean (infinite) prime in the Arakelov intersection pairing. Together with the N\'eron--Tate height (positivity at finite primes, cf.\ the companion BSD paper), this suggests a unified arithmetic positivity principle.
\end{remark}



\subsection{Directions for stronger positivity}
\label{sec:stronger_positivity}

Since $\mathrm{AP}^p = \mathrm{AbsHodge}^p$ (Remark~\ref{rem:ap_equals_absolute}), the current AP conditions do not go beyond Deligne's absoluteness. This section surveys candidate \textit{strengthening} directions that could produce genuinely new constraints beyond the Hodge-Riemann relations.

\subsubsection{\texorpdfstring{$p$-adic slope constraints}{p-adic slope constraints}}

\begin{remark}[$p$-adic positivity: Newton versus Hodge polygon]
\label{rem:newton-hodge}
Let $X/\mathbb{Q}$ be a smooth projective variety with good reduction at a prime $p$. The crystalline cohomology $H^{2p}_{\mathrm{cris}}(X_p/\mathbb{Z}_p)$ carries a Frobenius-linear endomorphism $\varphi$, whose eigenvalue slopes are constrained by the Newton polygon. The Katz conjecture (proved by Mazur~\cite{Mazur1972} and Ogus) states that the Newton polygon lies on or above the Hodge polygon.

Define a \textit{$p$-adic slope constraint} (AP$'$): a rational $(p,p)$-class $\alpha$ is \textit{$p$-adically positive} if, for all primes $\ell$ of good reduction, the $\ell$-adic realization $\alpha_\ell$ lies in the ``algebraic slope range'' -- i.e., its crystalline Frobenius slope equals $p$ (an integer). Algebraic classes automatically satisfy this (the Frobenius slope of an algebraic cycle of codimension $p$ is exactly $p$). For non-algebraic absolute Hodge classes, the slope could deviate from $p$, providing a new obstruction (O4) beyond absoluteness.

The challenge is that crystalline positivity is inherently $p$-adic and does not directly interact with the archimedean Hodge-Riemann form. An adelic formulation combining (AP1)--(AP3) at the archimedean place with slope constraints at finite places would be needed.
\end{remark}

\subsubsection{Arakelov positivity}

\begin{remark}[Arakelov positivity as global amplifier]
\label{rem:arakelov-positivity}
The Arakelov intersection pairing on arithmetic varieties combines:
\begin{itemize}
\item \textit{Finite primes:} intersection multiplicities in the fibres of a regular model $\mathcal{X} \to \mathrm{Spec}(\mathbb{Z})$;
\item \textit{Infinite prime:} the Green current at the archimedean place, related to $Q_L$ via the Hodge metric.
\end{itemize}
An \textit{Arakelov-strengthened AP condition} would require positivity of the \textit{global} Arakelov height pairing $\hat{h}(\alpha, \alpha) \geq 0$ for all models, not just the archimedean component. Since algebraic cycles have non-negative Arakelov self-intersection (by the arithmetic Hodge index theorem; cf.\ Faltings--Hriljac for surfaces, Moriwaki~\cite{Moriwaki2000} for the arithmetic height theory in higher dimensions), this is a natural candidate for AP$'$.

The difficulty is that the global pairing requires a regular model, which may not exist for arbitrary smooth projective varieties over $\mathbb{C}$ (descent to a number field is needed, and regularity at bad primes is not automatic).
\end{remark}

\begin{remark}[Test case for AP$'$]
\label{rem:test-case-ap-prime}
Any strengthened condition AP$'$ must satisfy two criteria:
\begin{enumerate}
\item \textit{Non-triviality:} AP$'$ must \textit{not} be automatically satisfied by every absolute Hodge class (otherwise it reduces to $\mathrm{AP}^p = \mathrm{AbsHodge}^p$ again).
\item \textit{Algebraic saturation:} Every algebraic class must satisfy AP$'$.
\end{enumerate}
A test case is needed: a variety $X$ with an absolute Hodge class $\alpha$ that is conjectured to be algebraic but whose algebraicity is not known. If AP$'$ could be \textit{verified} for $\alpha$ independently (e.g., via an explicit slope calculation), this would provide evidence for AP$' \Rightarrow$ algebraic. A candidate is the Hodge class on a general abelian fourfold ($g = 4$, $p = 2$) that lies outside the Lefschetz ring but is absolute by Deligne's theorem.
\end{remark}

\subsubsection{Tannakian perspective}

\begin{remark}[Tannakian reconstruction for CM cases]
\label{rem:tannakian}
For an abelian variety $A$ of CM type, the Mumford-Tate group $\mathrm{MT}(A)$ is a torus, and the category of Hodge structures generated by $H^1(A)$ is a neutral Tannakian category over $\mathbb{Q}$. The Hodge classes are precisely the $\mathrm{MT}(A)$-invariants.

The arithmetic positivity conditions can be reformulated in terms of the Mumford-Tate group: a class $\alpha$ is arithmetically positive if and only if it lies in the $\mathrm{MT}(A)$-invariant subspace and satisfies the positivity constraints at every embedding $\sigma: \mathbb{Q}(\mathrm{End}(A)) \hookrightarrow \mathbb{C}$ (cf.\ Deligne--Milne~\cite{DeligneMilne1982}). For CM abelian varieties, this gives an explicit criterion: $\alpha$ is AP if and only if it is fixed by the Taniyama group action -- which is exactly absoluteness.

To go beyond this, one would need \textit{motivic} positivity constraints that are not captured by the Mumford-Tate group alone. The conjectural category of mixed motives, equipped with its weight filtration and realizations, could provide such constraints via the positivity of the motivic height pairing (conjectured by Beilinson).
\end{remark}

\begin{remark}[Mumford-Tate analysis for Conjecture~\ref{conj:ghr}]
\label{rem:mumford-tate-ghr}
For abelian varieties $A$, Deligne's theorem gives
$H^{2p}(A,\mathbb{Q})\cap H^{p,p}(A)=\mathrm{AbsHodge}^p(A)=\mathrm{AP}^p(A)$.
Consequently the GHR obstruction cannot distinguish a hypothetical non-algebraic Hodge class on an abelian variety: if such a class exists, it is already absolute and its GHR spectrum is non-negative by the Hodge-Riemann relations. Thus the obstruction is vacuous on abelian varieties as a proof tool; the Hodge Conjecture for abelian varieties is exactly the remaining assertion that these absolute/AP classes are algebraic.

The remaining challenge is thus purely the \textit{algebraicity} of absolute Hodge classes. Progress depends on the Mumford-Tate conjecture: if $\mathrm{MT}(A)$ acts faithfully on $H^*(A, \mathbb{Q})$ and the $\mathrm{MT}(A)$-invariant classes are algebraic, then the Hodge conjecture for $A$ follows. This is known for abelian varieties of dimension $\leq 5$ (Moonen~\cite{Moonen1999}) and for abelian varieties of CM type (Abdulali~\cite{Abdulali1994}).
\end{remark}

\subsubsection{Prismatic and categorical approaches}

\begin{remark}[Prismatic positivity (AP4)]
\label{rem:prismatic}
Bhatt and Scholze's prismatic cohomology~\cite{BhattScholze2022} -- building on the integral $p$-adic Hodge theory of Bhatt--Morrow--Scholze~\cite{BhattMorrowScholze2018} -- provides a unified framework for $p$-adic cohomology theories (crystalline, de Rham, \'etale). In this framework, a Hodge class $\alpha$ has a prismatic realization $\alpha_{\Delta}$ in $H^{2p}_{\Delta}(X/\mathbb{Z}_p)$.

A conjectural \textit{prismatic positivity condition} (AP4) would require: for every prime $p$ and every prism $(A, I)$, the prismatic realization $\alpha_{\Delta}$ lies in a ``positive cone'' defined by the Nygaard filtration and the Frobenius action. Algebraic classes satisfy this automatically (they lie in the image of the prismatic cycle class map). Non-algebraic absolute Hodge classes could potentially violate AP4, providing a new obstruction beyond (O1)--(O3).

This is speculative but motivated by the fact that prismatic cohomology ``sees'' arithmetic structure invisible to classical Hodge theory. The compatibility with (AP1)--(AP3) at the archimedean place would need to be established via the prismatic--de Rham comparison.
\end{remark}

\begin{remark}[Prismatic cohomology route]\label{rem:prismatic-route}
Nie \cite{Nie2021} constructed absolute Hodge cycles in the prismatic cohomology of abelian varieties over $p$-adic fields, using Andr\'e's motivated cycles within the Bhatt--Scholze framework. Prismatic cohomology $R\Gamma_\Delta(X/S)$ interpolates \'etale, de Rham, and crystalline cohomology integrally, providing a universal cohomological framework. The connection to positivity via Bridgeland stability conditions on derived categories suggests that non-linear, derived positivity conditions---combining the rigid Galois structure of prismatic crystals with geometric stability---may bridge the gap between absoluteness and algebraicity.
\end{remark}

\subsubsection*{Threshold desideratum: beyond Deligne}

\begin{remark}[Desideratum: Language Bridge -- JD-Hodge]\label{ax:jd-hodge}
Ideally, one would like a category $\mathcal{C}$ of ``Hodge carriers'' (objects representing cohomological data), equipped with:
\begin{enumerate}
\item[(JD1)] \textbf{Functoriality:} Pullback, pushforward, and products act on $\mathcal{C}$ without non-canonical choices.
\item[(JD2)] \textbf{Intrinsic filtration:} A functorial Hodge filtration $F^\bullet$ is defined canonically (not via a choice of embedding or model).
\item[(JD3)] \textbf{Canonical polarisation:} A positivity structure compatible with (JD1) and (JD2) exists, generalising the Hodge--Riemann form.
\end{enumerate}
In such a framework, the Hodge conjecture would become a \emph{formal consequence}: positivity in $\mathcal{C}$ would imply algebraicity. Deligne's absoluteness would then be a shadow (necessary condition) of the richer structure in $\mathcal{C}$.
\end{remark}

\begin{remark}[The non-canonicity scanner]\label{rem:noncanonicity-scanner}
To locate where the current theory falls short of JD-Hodge, scan for three signatures:
\begin{enumerate}
\item \textbf{Choice dependence:} Where do splittings or comparison isomorphisms (de Rham/Betti/$\ell$-adic) appear that ``exist'' but are not natural transformations?
\item \textbf{Incompatible operations:} Which geometric operation (pullback under non-finite morphisms, Gysin pushforward) destroys the polarisation or destabilises the filtration?
\item \textbf{Comparison breaks:} Where does the statement depend on a comparison between realisations that is not structurally controlled?
\end{enumerate}
In the current AP framework, the break occurs at AP2a (absoluteness): the comparison isomorphism $H^*_{\mathrm{dR}}(X) \otimes \mathbb{C} \cong H^*_{\mathrm{Betti}}(X^\sigma) \otimes \mathbb{C}$ is \emph{not} a morphism in any natural category---it requires the choice of $\sigma \in \mathrm{Aut}(\mathbb{C}/\mathbb{Q})$.
\end{remark}

\begin{remark}[Candidates for $\mathcal{C}$]\label{rem:candidates-C}
Three candidates for the category $\mathcal{C}$ are currently under investigation:
\begin{enumerate}
\item \textbf{Prismatic cohomology} (Bhatt--Scholze, Nie \cite{Nie2021}): $R\Gamma_\Delta(X/S)$ interpolates all classical realisations integrally, with a canonical Frobenius action providing (JD2). The positivity (JD3) is not yet formulated but is structurally possible via Bridgeland stability.
\item \textbf{Motivic cohomology} (Voevodsky): The derived category of motives $\mathrm{DM}(k)$ satisfies (JD1) by construction. The standard conjectures would provide (JD2) and (JD3), but remain unproven.
\item \textbf{Mixed Hodge modules} (Saito): Satisfy (JD1)--(JD2) for smooth projective varieties, but (JD3) is limited to the classical HR form---precisely the AP$=$AbsHodge barrier of this paper.
\end{enumerate}
\end{remark}

\begin{lemma}[JD-Hodge implies the Hodge conjecture for abelian varieties]\label{lem:jd-abelian}
If $\mathcal{C}$ satisfies (JD1)--(JD3) and is compatible with the Tannakian structure of the category of abelian motives, then every Hodge class on an abelian variety is algebraic.
\end{lemma}

\begin{proof}[Proof sketch]
Let $A$ be an abelian variety and $\alpha \in H^{2p}(A, \mathbb{Q}) \cap H^{p,p}(A)$ a Hodge class. By Deligne's theorem (Theorem~\ref{thm:deligne_absolute}), $\alpha$ is absolute, hence lies in $\mathrm{AP}^p(A)$ by Proposition~\ref{prop:absolute_implies_ap}. In the category $\mathcal{C}$, the canonical polarisation (JD3) and functoriality (JD1) together with the Tannakian formalism imply that $\alpha$ is a morphism in $\mathcal{C}$, i.e., motivic. By the Tannakian reconstruction for the category of abelian motives (cf.\ Deligne--Milne~\cite{DeligneMilne1982}), motivic classes on abelian varieties are algebraic.
\end{proof}

\begin{remark}[The Deligne boundary as a language defect]\label{rem:deligne-language}
The JD-Hodge desideratum (Remark~\ref{ax:jd-hodge}) identifies the obstruction as \emph{linguistic}, not computational: the comparison isomorphisms $\mathrm{comp}_{\mathrm{dR} \leftrightarrow B}$ and $\mathrm{comp}_{\ell\text{-adic} \leftrightarrow B}$ exist but are not functorial in $X$. Absolute Hodge classes are precisely those compatible under all realisations---and algebraicity becomes $\mathrm{AH}(X) = \mathrm{Im}(\mathrm{CH}^*(X) \to H^*(X))$. If the comparisons were natural transformations in a category $\mathcal{C}_{\mathrm{abs}}$, this would be a formal consequence. The Deligne boundary is therefore: \emph{non-canonical splittings make every algebraicity proof unstable under functors}.
\end{remark}

\begin{remark}[Bridgeland stability and positivity on derived categories]
\label{rem:bridgeland}
Bridgeland stability conditions~\cite{Bridgeland2007} on the derived category $D^b(X)$ provide a notion of ``positivity'' for objects (sheaves, complexes) rather than cohomology classes. The central charge $Z: K_0(X) \to \mathbb{C}$ maps the Grothendieck group to $\mathbb{C}$, and stability requires $\mathrm{Im}(Z) > 0$ or $\mathrm{Im}(Z) = 0, \mathrm{Re}(Z) < 0$ for semistable objects.

A potential strengthening of AP: require that the cohomology class $\alpha$ arises as the Chern character of a Bridgeland-stable object (or a linear combination thereof). This is strictly stronger than being a Hodge class, because not every Hodge class is the Chern character of a stable object. However, connecting this to the AP framework requires understanding the image of the Chern character map in the presence of stability conditions -- a problem related to the ``Bogomolov-type inequalities'' in higher dimensions.
\end{remark}


\subsubsection{Geometric and o-minimal approaches}

\begin{remark}[VHS families and monodromy]
\label{rem:vhs-monodromy}
Let $f: \mathcal{X} \to S$ be a smooth projective family over a quasi-projective base $S$, giving rise to a variation of Hodge structure (VHS). A Hodge class $\alpha_s \in H^{2p}(X_s, \mathbb{Q}) \cap H^{p,p}(X_s)$ at a point $s \in S$ is \textit{globally flat} if it extends to a flat section of the local system $R^{2p}f_*\mathbb{Q}$ over a neighborhood of $s$, and \textit{monodromically stable} if it is invariant under the monodromy representation $\pi_1(S, s) \to \mathrm{GL}(H^{2p}(X_s, \mathbb{Q}))$.

The theorem of Cattani--Deligne--Kaplan~\cite{CattaniDeligneKaplan1995} shows that the locus of algebraic classes is a countable union of closed algebraic subvarieties of $S$. A strengthened AP condition could require that $\alpha_s$ extends as a \textit{flat section} over a Zariski-dense subset of $S$ (not merely that it is a Hodge class at $s$). This is related to the semisimplicity of the monodromy representation and the algebraicity of flat sections predicted by the Hodge conjecture for families.
\end{remark}

\begin{remark}[$p$-adic positivity via o-minimality]
\label{rem:o-minimality}
The theory of o-minimal structures (Bakker--Brunebarbe--Tsimerman~\cite{BBT2023}) has been applied to prove definability results for period maps and Hodge loci. In this framework, the period domain is a definable analytic space, and the Hodge locus is a definable subset. The key insight is that definable sets in o-minimal structures have strong tameness properties (e.g., finite number of connected components, controlled topology).

A potential AP$'$ condition: require that the ``orbit'' of $\alpha$ under the definable period map is algebraic (i.e., the Zariski closure of the set $\{s \in S : \alpha_s \text{ is Hodge}\}$ is an algebraic subvariety). By CDK, this is true for algebraic classes. For non-algebraic Hodge classes, the orbit could potentially be ``wild'' (transcendental), violating the tameness condition. This connects to Binyamini--Novikov's work~\cite{BinyaminiNovikov2017} on counting rational points on definable sets.
\end{remark}

\begin{remark}[Galois defects for transcendental classes]
\label{rem:galois-defects}
If a Hodge class $\alpha$ is not absolute, then there exists $\sigma \in \mathrm{Aut}(\mathbb{C})$ such that $\alpha^\sigma$ is not of type $(p,p)$ on $X^\sigma$. The ``size'' of this defect -- measured, for instance, by the component of $\alpha^\sigma$ in $H^{p-1,p+1}(X^\sigma) \oplus H^{p+1,p-1}(X^\sigma)$ -- is a quantitative measure of non-algebraicity. If this defect grows under iteration of Galois conjugation (i.e., the orbit $\{\alpha^{\sigma^n}\}$ generates a ``macroscopic'' subspace of $H^{2p}(X, \mathbb{C})$), then the GHR spectrum would become negative, providing a computable obstruction.

This connects to the question of whether transcendental Hodge classes can be ``isolated'' (small Galois orbit) or must generate large representations of $\mathrm{Aut}(\mathbb{C})$. The latter would make the GHR obstruction effective for transcendental classes.
\end{remark}


\subsection{Beyond bilinear positivity: nonlinear and averaged conditions}
\label{sec:nonlinear-positivity}

The Hodge--Riemann relations provide \emph{bilinear} positivity: $(-1)^p Q_L(\alpha, \bar\alpha) > 0$ for primitive $(p,q)$-classes. As noted in the limitations (Section~9), this does not suffice for the full Hodge conjecture. We identify three directions for \emph{nonlinear} or \emph{averaged} positivity that go beyond HR:

\begin{remark}[Three routes to nonlinear positivity]
\label{rem:nonlinear-positivity}
\emph{Boundary positivity.}
For a family of ample classes $L_t \to L_0$, where $L_0$ is nef but not ample, the signed values $(-1)^p Q_{L_t}$ can have a non-negative limit as $t \to 0$ even when $Q_{L_0}$ is not defined. This captures classes that are algebraic in the limit but not detected by any single polarisation.

\emph{Positivity under motivic convolution.}
For a product variety $X \times Y$, the K\"unneth decomposition gives $H^{p,p}(X \times Y) = \bigoplus_{\substack{a+c=p \\ b+d=p}} H^{a,b}(X) \otimes H^{c,d}(Y)$. For \emph{rational} $(p,p)$-classes restricted to the diagonal summands $H^{a,a}(X) \otimes H^{p-a,p-a}(Y)$, the positivity condition on the product is multiplicative (Lemma~\ref{lem:sign_product}). The cross terms $H^{a,b}(X) \otimes H^{c,d}(Y)$ with $a \neq b$ contribute to the $(p,p)$-type only when the mixed Hodge types cancel, and their positivity analysis requires going beyond the diagonal framework of this paper.

\emph{Positivity after Galois averaging.}
For an absolutely Hodge class $\alpha$, a suitably defined finite-level Galois average $\bar\alpha$ (cf.\ Remark~\ref{rem:galois-profinite}) is again an absolute Hodge class. For $p = 1$, such averaged classes are algebraic by the Lefschetz $(1,1)$-theorem. For $p \geq 2$, the algebraicity of the average is \emph{not known} in general---it is equivalent to (a special case of) Deligne's question. The arithmetic positivity conditions AP1--AP3 can be reformulated as a proximity heuristic to this average in the $Q_L$-metric, but the reformulation is useful only if the average can be shown to be algebraic by an argument independent of the Hodge conjecture.
\end{remark}

\begin{remark}[Galois averaging: profinite limit]\label{rem:galois-profinite}
The na\"ive averaging $\bar{\alpha} = \frac{1}{|\mathrm{Aut}(\mathbb{C}/\mathbb{Q})|} \sum_\sigma \sigma^*\alpha$ is ill-defined since $\mathrm{Aut}(\mathbb{C}/\mathbb{Q})$ is a profinite group of cardinality $2^{\mathfrak{c}}$. A cautious formulation works only at finite level: for each finite Galois extension $K/\mathbb{Q}$ over which the relevant model and comparison data are controlled, form
\[
\bar{\alpha}_K = \frac{1}{|\mathrm{Gal}(K/\mathbb{Q})|}
  \sum_{\sigma \in \mathrm{Gal}(K/\mathbb{Q})} \sigma^*\alpha.
\]
Any projective-limit interpretation requires additional descent and compatibility data; it should not be treated as an automatic algebraic class.
\end{remark}


\subsection{Numerical illustration}

Since $\mathrm{AP}^p = \mathrm{AbsHodge}^p$ is a \emph{theorem} (not a conjecture), the following numerical computations serve as \textbf{pedagogical illustrations}, not as empirical verification. They demonstrate the AP framework concretely and illustrate the GHR spectrum on specific varieties.

The Python script \texttt{compute\_ghr\_spectrum.py} (accompanying this paper) computes the AP conditions across selected test cases:
\begin{itemize}
\item K3 surfaces (4 polarizations), abelian surfaces (5 polarizations), CM endomorphism classes on $E \times E$;
\item $K3 \times K3$ tensor products with non-tensor-product perturbations.
\end{itemize}

All test configurations confirm $(-1)^p Q_L \geq 0$ on $P^{p,p}$, as guaranteed by the Hodge-Riemann relations. The $K3 \times K3$ test with non-tensor-product perturbations -- classes that are \emph{not} of type $(p,p)$ after $\sigma$-conjugation, and hence not absolute Hodge -- produces negative eigenvalues of $(-1)^p Q_L$, illustrating that the GHR spectrum can detect non-absoluteness. This does not go beyond Deligne's absoluteness criterion: the negative eigenvalues arise precisely because the perturbed classes fail to be $(p,p)$-classes on the conjugate variety, so the Hodge-Riemann relations do not apply.

\subsection{Controlled test-family transfer beyond \texorpdfstring{$\mathrm{AP}=\mathrm{AbsHodge}$}{AP=AbsHodge}}
\label{subsec:zookeeper-hodge-transfer}

The post-Zookeeper interpretation of this paper is deliberately modest.
Since arithmetic positivity collapses to absolute Hodge classes, the next
step is not another global positivity axiom, but a test-family ladder for
candidate strengthenings beyond Deligne's absoluteness.
The symbols $\mathrm{AP}'$, $\mathrm{AP}''$, and
$\mathrm{AP}^{\prime\prime\prime}$ do not denote proposed global
replacements for the Hodge conjecture. They denote increasingly restrictive
tests on controlled families: $\mathrm{AP}'$ records provenance defects,
$\mathrm{AP}''$ asks for a bounded orbit-rigidity certificate, and
$\mathrm{AP}^{\prime\prime\prime}$ filters out certificates whose auxiliary
choices merely encode the target class. None of these conditions is claimed
to imply the Hodge conjecture in general.

Thus the notation is a claim-status convention, not a theorem label. Any
algebraicity conclusion in this subsection is only as strong as the named
family hypotheses that are verified independently of the target class
$\alpha$. If the auxiliary data are selected after seeing $\alpha$, the
condition is treated as a failed test rather than as evidence.

The current CORE closure of Riemann-$C2/MS2$ does not strengthen this
paper's Hodge result.  Hodge remains a negative control for the programme:
Hodge--Riemann positivity saturates too broadly, so any future transfer
must add genuinely new arithmetic or motivic constraints rather than reuse
the Riemann kernel closure.

The five transfer slots are as follows.
\begin{description}
  \item[Operator/form.] The Hodge Laplacian, Mumford--Tate and monodromy
  operators, and the period map.
  \item[Defect.] The residual transcendental part of an absolute Hodge
  class after subtracting the algebraic cycle span known in the family.
  \item[Approximation.] Abelian fourfolds of Weil type, singular
  OG6-type hyperkaehler models, powers of known examples, and reductions
  in characteristic $p$.
  \item[Gap.] Extra constraints beyond Hodge--Riemann positivity:
  $p$-adic slopes, Tate constraints, Mumford--Tate restrictions, and
  o-minimal algebraicity.
  \item[Limit.] Deformation, powers, and specialization from a controlled
  family to a broader motivic setting.
\end{description}

Accordingly, a useful next version should contain a table of test
families.  For each family it should record: known Hodge/Tate status, the
cycle mechanism used in the literature, the proposed strengthened
condition $\mathrm{AP}'$, and whether the defect vanishes.  Weil fourfolds
of discriminant one, OG6-type constructions, abelian varieties over finite
fields satisfying the standard conjecture of Hodge type, and Voisin-type
stress tests form the natural first ladder.  If $\mathrm{AP}'$ still holds
automatically on all negative controls, it is too weak; if it separates a
known positive family from a stress case, it becomes a genuine candidate
for further formalisation.

\subsubsection*{Towards \texorpdfstring{$\mathrm{AP}'$}{AP'}: Provenance defects}
\label{subsec:provenance-defects}

The first strengthening should be read as a provenance scanner rather than
as another positivity cone. Relative to a named test family, define a defect
vector
\[
D_{\mathrm{prov}}(\alpha)
  = (D_{\mathrm{slope}},D_{\mathrm{Tate}},D_{\mathrm{MT}},
     D_{\mathrm{mon}},D_{\mathrm{cycle}}),
\]
where the entries record, respectively, slope compatibility at good
reductions, explanation by the Tate/Galois fixed part, Mumford--Tate
compatibility, monodromy or o-minimal persistence, and the residual class
left after subtracting the known algebraic cycle span.  A conditional
fingerprint statement would say that, in a test family where the relevant
Tate/standard-conjecture input and comparison functoriality are independently
available, $D_{\mathrm{prov}}(\alpha)=0$ places $\alpha$ in the known
algebraic cycle span.  This is a special-family diagnostic, not a global
theorem.

Each component of $D_{\mathrm{prov}}$ must be recorded with status
\emph{verified}, \emph{failed}, or \emph{unknown}.  Unknown components do not
count as zero defects.  This convention prevents the provenance scanner from
becoming tautological: $\mathrm{AP}'$ is passed only when all required
components are independently verified in the chosen test family.

\subsubsection*{Towards \texorpdfstring{$\mathrm{AP}''$}{AP''}: Conditional Orbit-Rigidity}
\label{subsec:orbit-rigidity}

The five transfer slots suggest a prototype sharpening.  Let
$f\colon \mathcal{X}_{S}\to S$ be a spread-out of $X$ over a smooth
base~$S$, let $Z\subset S$ be an irreducible algebraic component of the
Hodge locus of~$\alpha$, let $B$ denote the corresponding relative
Chow or Hilbert component, and let $T\subset Z$ be a Zariski-dense
subset at whose points $\alpha_{t}$ is algebraic relative to~$B$.
We call $\alpha$ \emph{orbit-rigid on $(f,Z,B)$} if:
\begin{enumerate}
\item[(i)] \emph{Properness.} The relative Chow/Hilbert components
  parametrising the relevant cycle classes are proper over~$S$.
\item[(ii)] \emph{Density.} $T$ is Zariski-dense in~$Z$.
\item[(iii)] \emph{Comparison functoriality.}  The comparison
  isomorphisms between Betti, de~Rham, $\ell$-adic, and crystalline
  realisations are functorial on the relevant quotient; no new residual
  classes arise when lifting from a special to the generic fibre.
\item[(iv)] \emph{Specialisation rigidity.}  Base change along
  $Z\to S$ does not create new residual cycle classes.
\end{enumerate}

\begin{lemma}[Conditional Orbit-Rigidity]
\label{lem:orbit-rigidity}
Let $\alpha\in\mathrm{AbsHodge}^{p}(X)$ be orbit-rigid on $(f,Z,B)$
in the sense above.  Then the restriction $\alpha_{\eta}$ to the generic
fibre $\mathcal{X}_{\eta}$ over the generic point $\eta$ of~$Z$ lies in
the image of the cycle class map
$\mathrm{cl}\colon\mathrm{CH}^{p}(\mathcal{X}_{\eta})_{\mathbb{Q}}\to
H^{2p}(\mathcal{X}_{\eta},\mathbb{Q}(p))$.
If the base-point $0\in S$ lies in the same controlled component,
then $\alpha$ is algebraic.
\end{lemma}

\begin{proof}[Proof idea]
Properness of $B$ \emph{(i)} and Zariski-density of $T$ \emph{(ii)} imply via
the valuative criterion that the closure of the algebraic locus in $B$
contains the generic point~$\eta$ of~$Z$.  Comparison functoriality \emph{(iii)}
ensures that each algebraic class over $t\in T$ lifts consistently in all
cohomological realisations.  Specialisation rigidity \emph{(iv)} prevents
accumulation of transcendental residual classes at the limit.  A complete
proof requires the Cattani--Deligne--Kaplan principle for Hodge loci
\cite{CattaniDeligneKaplan1995} together with a properness argument for
the relative Hilbert scheme; both
steps are conditional on the standard conjectures in the relevant
characteristic-$p$ reductions.
\end{proof}

\begin{remark}
\label{rem:ap-double-prime}
Lemma~\ref{lem:orbit-rigidity} does not prove the Hodge Conjecture.
It should be read as a conditional closure template, not as a new
unconditional algebraicity theorem. Its force lies entirely in the
independent verification of bounded relative Chow/Hilbert data, comparison
functoriality, and specialisation rigidity in a named test family. Under
those independently checked hypotheses, algebraicity propagates from a
Zariski-dense set to the generic point.  Weil abelian fourfolds of
discriminant one are a natural first candidate for a positive test;
$K3\times K3$ perturbations outside the tensor-product locus serve as the
first negative control.  Whether AP$''$ separates these two families
determines whether it constitutes a genuine strengthening beyond
$\mathrm{AbsHodge}$.
\end{remark}

\subsubsection*{Towards \texorpdfstring{$\mathrm{AP}^{\prime\prime\prime}$}{AP'''}: Uniformity and anti-advice}
\label{subsec:ap-triple-prime}

The orbit-rigidity certificate still has a failure mode: the family
$f$, the Hodge-locus component $Z$, the Chow/Hilbert component $B$, or the
dense set $T$ could be chosen after seeing the target class $\alpha$.  In
that case the certificate may encode $\alpha$ rather than explain it.  We
therefore use
\[
\mathrm{AP}^{\prime\prime\prime}_{\mathrm{unif}}(\alpha;R,C)
\]
as an anti-tautology filter for an $\mathrm{AP}''$ certificate $C$. It
requires that the auxiliary data are generated from a pre-registered rule
class $R$ depending only on the ambient data $(X,p)$ or on a fixed
test-family ladder, with bounded advice, erosion stability under removal of
an $\alpha$-specific incidence condition, support on a neighbouring family
or Zariski-dense set, and failure on negative controls driven only by
Hodge--Riemann form stability.  This condition does not add a new global
Hodge conjecture; it says when a special-family certificate is non-tautological.

Operationally, a certificate fails the $\mathrm{AP}^{\prime\prime\prime}$
filter if any of the following checks fails:
\begin{enumerate}
\item \emph{Pre-registration:} the rule class $R$ was not fixed before the
  target class or does not follow from ambient data.
\item \emph{Bounded advice:} the choices of $f,Z,B,T$ carry enough
  information to reconstruct $\alpha$.
\item \emph{Erosion stability:} removing an $\alpha$-specific incidence
  condition destroys the certificate.
\item \emph{Neighbour-swarm:} the same rule does not support a neighbouring
  family or a Zariski-dense set.
\item \emph{Negative controls:} Voisin/$K3\times K3$-type stress families
  pass merely because of Hodge-Riemann form stability.
\end{enumerate}

\begin{center}
\resizebox{\columnwidth}{!}{%
\begin{tabular}{llllll}
\toprule
\textbf{Family} & \textbf{Cycle mechanism} & \textbf{AP$'$} & \textbf{AP$''$} & \textbf{AP$^{\prime\prime\prime}$} & \textbf{Status} \\
\midrule
CM/Weil fourfolds & endomorphism cycles & slope/Tate/MT & candidate & rule from endomorphisms & positive control \\
Abelian fourfolds & known cycle span & Tate/MT & candidate & ambient family rule & positive control \\
OG6-type families & hyperkaehler cycles & monodromy/cycle & open & needs registry & test case \\
Char-$p$ Tate cases & Tate classes & slope/Tate & conditional & reduction rule & arithmetic test \\
Voisin/$K3\times K3$ stress & none fixed & residual class & should fail & post-hoc risk & negative control \\
\bottomrule
\end{tabular}}
\end{center}

\subsection*{Result classification}

This paper is a \textbf{boundary result}: it determines precisely what positivity-based methods can and cannot achieve for the Hodge conjecture.

\begin{center}
\resizebox{\columnwidth}{!}{%
\begin{tabular}{lll}
\toprule
\textbf{Result} & \textbf{Type} & \textbf{Use in this paper} \\
\midrule
\multicolumn{3}{l}{\textit{No-go / boundary result:}} \\
AP $=$ AbsHodge (HR adds nothing) & Theorem & proved boundary \\
HC reduces to Deligne's question & Corollary & scope limiter \\
\midrule
\multicolumn{3}{l}{\textit{Constructive tools:}} \\
GHR spectrum definition & Definition & visualization only \\
Finite pullback / primitive K\"unneth & Partial, conditional & limited tests \\
Voisin-type GHR detection & Numerical & illustration, not criterion \\
Conditional Orbit-Rigidity $(\mathrm{AP}'')$ & Lemma (conditional) & template \\
Uniform anti-advice filter $(\mathrm{AP}''')$ & Methodological criterion & admissibility filter \\
\midrule
\multicolumn{3}{l}{\textit{Open (equivalent to Deligne 1982):}} \\
AbsHodge $\stackrel{?}{=}$ Algebraic & Open & not addressed \\
Stronger positivity beyond HR & Research direction & future work \\
\bottomrule
\end{tabular}}
\end{center}


% ===================================================================
% 10. CONCLUSION
% ===================================================================
\section{Conclusion}
\label{sec:conclusion}

We have proposed a reformulation of the Hodge Conjecture through the lens of positivity and No-Go obstructions. The key concept -- \textit{arithmetic positivity} -- combines Hodge-Riemann positivity, absolute stability under $\mathrm{Aut}(\mathbb{C})$, and Lefschetz compatibility into a single condition.

A central finding of this paper (Remark~\ref{rem:ap_equals_absolute}) is that the arithmetically positive cone coincides with Deligne's absolute Hodge classes: $\mathrm{AP}^p(X) = \mathrm{AbsHodge}^p(X)$. This means the Arithmetic Positivity Conjecture is equivalent to Deligne's question ``Are all absolute Hodge classes algebraic?'' The positivity conditions (AP1), (AP2b), and (AP3) are automatically satisfied by any absolute Hodge class, via the Hodge-Riemann bilinear relations.

The paper does not introduce a genuinely new conjecture; the saturation of Hodge-Riemann positivity for $(p,p)$-classes is a standard consequence of the HR relations. The paper's contributions are expository and systematic:

\begin{enumerate}
\item \textbf{A new perspective:} The No-Go formulation recasts the Hodge Conjecture from a construction problem to an impossibility statement.

\item \textbf{A visualization of the absoluteness test:} The Galois-Hodge-Riemann spectrum reformulates Deligne's absoluteness criterion in the language of Hodge-Riemann positivity, making explicit which Galois conjugation could witness non-algebraicity. It does not provide constraints beyond Deligne's absoluteness.

\item \textbf{Partial functoriality checks:} The paper records partial functoriality checks rather than a full functorial theory: finite pullback is conditional on the ample-cone comparison, primitive K\"unneth products are controlled for product polarizations, and finite pushforward remains open.
\end{enumerate}

The remaining challenge is Deligne's question: proving that absolute Hodge classes are algebraic. Future work should seek stronger positivity conditions -- beyond the Hodge-Riemann relations -- that could further narrow the gap between absolute Hodge classes and algebraic classes.

To summarize the philosophical position of this paper: the Hodge Conjecture, viewed through the positivity lens, does not ask us to build algebraic cycles -- it asks us to show that geometry has no room for anything else. The present work shows that Hodge-Riemann positivity alone cannot accomplish this; genuinely new positivity conditions are required.


% ===================================================================
% ACKNOWLEDGMENTS
% ===================================================================
\begin{acknowledgments}
This work was conducted as independent research without institutional funding.

\textit{AI tools.} AI-based tools (Claude by Anthropic) were used for mathematical formalization and text generation. All mathematical content, conjectures, and responsibility for correctness lie solely with the author.

\textit{Funding information.} This research received no external funding.
\end{acknowledgments}


% ===================================================================
% BIBLIOGRAPHY
% ===================================================================
\begin{thebibliography}{30}

\bibitem{Hodge1950}
W.~V.~D.~Hodge (1950).
The topological invariants of algebraic varieties.
\textit{Proceedings of the International Congress of Mathematicians (Cambridge, MA, 1950)}, vol.~1, pp.\ 181--192. American Mathematical Society.

\bibitem{Voisin2002}
C.~Voisin (2002).
\textit{Hodge Theory and Complex Algebraic Geometry I, II}.
Cambridge University Press.
DOI: 10.1017/CBO9780511615344.

\bibitem{GH1978}
P.~Griffiths, J.~Harris (1978).
\textit{Principles of Algebraic Geometry}.
John Wiley \& Sons, New York.

\bibitem{Deligne1982}
P.~Deligne (1982).
Hodge cycles on abelian varieties (Notes by J.~S.~Milne).
In \textit{Hodge Cycles, Motives, and Shimura Varieties}, Lecture Notes in Mathematics 900, pp.\ 9--100. Springer.
DOI: 10.1007/978-3-540-38955-2\_3.

\bibitem{Andre1996}
Y.~Andr\'e (1996).
Pour une th\'eorie inconditionnelle des motifs.
\textit{Publications Math\'ematiques de l'IH\'ES}, 83, 5--49.
DOI: 10.1007/BF02698643.

\bibitem{Charles2013}
F.~Charles (2013).
The Tate conjecture for K3 surfaces over finite fields.
\textit{Inventiones Mathematicae}, 194(1), 119--145.
DOI: 10.1007/s00222-012-0443-y.

\bibitem{Faltings1983}
G.~Faltings (1983).
Endlichkeitss\"atze f\"ur abelsche Variet\"aten \"uber Zahlk\"orpern.
\textit{Inventiones Mathematicae}, 73, 349--366.
DOI: 10.1007/BF01388432.

\bibitem{Abdulali1994}
S.~Abdulali (1994).
Algebraic cycles in families of abelian varieties.
\textit{Canadian Journal of Mathematics}, 46(6), 1121--1134.
DOI: 10.4153/CJM-1994-063-0.

\bibitem{Hazama1994}
F.~Hazama (1994).
The generalized Hodge conjecture for stably nondegenerate abelian varieties.
\textit{Compositio Mathematica}, 93(2), 129--137.

\bibitem{Moonen1999}
B.~Moonen (1999).
Notes on Mumford-Tate groups.
Preprint, University of Amsterdam.

\bibitem{Voisin2002b}
C.~Voisin (2002).
A counterexample to the Hodge conjecture extended to K\"ahler varieties.
\textit{International Mathematics Research Notices}, 2002(20), 1057--1075.
DOI: 10.1155/S1073792802111135.

\bibitem{CattaniDeligneKaplan1995}
E.~Cattani, P.~Deligne \& S.~Kaplan (1995).
On the locus of Hodge classes.
\textit{Journal of the American Mathematical Society}, 8(2), 483--506.
DOI: 10.2307/2152824.

\bibitem{Grothendieck1969}
A.~Grothendieck (1969).
Standard conjectures on algebraic cycles.
In \textit{Algebraic Geometry (Bombay, 1968)}, pp.\ 193--199. Oxford University Press.

\bibitem{GrossZagier1986}
B.~Gross \& D.~Zagier (1986).
Heegner points and derivatives of $L$-series.
\textit{Inventiones Mathematicae}, 84, 225--320.
DOI: 10.1007/BF01388809.

\bibitem{Deligne1974}
P.~Deligne (1974).
La conjecture de Weil: I.
\textit{Publications Math\'ematiques de l'IH\'ES}, 43, 273--307.
DOI: 10.1007/BF02684373.

\bibitem{FST-Unified2026}
L.~Geiger (2026).
FST Spectrum Duality: The Renormalized Free Energy Principle as Physical Instantiation of Functional Stability Theory.
Zenodo preprint, version 1.7.
\href{https://doi.org/10.5281/zenodo.19036190}{doi:10.5281/zenodo.19036190}.

\bibitem{GeigerRH2026c}
L.~Geiger (2026).
From Landscape to Proof: The Even-Dominance Route to the Riemann Hypothesis.
Zenodo preprint, version 2.3.
\href{https://doi.org/10.5281/zenodo.19035640}{doi:10.5281/zenodo.19035640}.

\bibitem{Mazur1972}
B.~Mazur (1972).
Frobenius and the Hodge filtration.
\textit{Bulletin of the American Mathematical Society}, 78(5), 653--667.
DOI: 10.1090/S0002-9904-1972-12976-8.

\bibitem{Moriwaki2000}
A.~Moriwaki (2000).
Arithmetic height functions over finitely generated fields.
\textit{Inventiones Mathematicae}, 140(1), 101--142.
DOI: 10.1007/s002220050358.

\bibitem{DeligneMilne1982}
P.~Deligne \& J.~S.~Milne (1982).
Tannakian categories.
In \textit{Hodge Cycles, Motives, and Shimura Varieties}, Lecture Notes in Mathematics 900, pp.\ 101--228. Springer.
DOI: 10.1007/978-3-540-38955-2\_4.

\bibitem{BhattScholze2022}
B.~Bhatt \& P.~Scholze (2022).
Prisms and prismatic cohomology.
\textit{Annals of Mathematics}, 196(3), 1135--1275.
DOI: 10.4007/annals.2022.196.3.5.

\bibitem{Bridgeland2007}
T.~Bridgeland (2007).
Stability conditions on triangulated categories.
\textit{Annals of Mathematics}, 166(2), 317--345.
DOI: 10.4007/annals.2007.166.317.

\bibitem{BBT2023}
B.~Bakker, B.~Brunebarbe \& J.~Tsimerman (2023).
o-minimal GAGA and a conjecture of Griffiths.
\textit{Inventiones Mathematicae}, 232(1), 163--228.
DOI: 10.1007/s00222-022-01166-1.

\bibitem{Nie2021} T.~Nie, \textit{Absolute Hodge Cycles in Prismatic Cohomology}, PhD Thesis, Harvard University, 2021.

\bibitem{BinyaminiNovikov2017}
G.~Binyamini \& D.~Novikov (2017).
Wilkie's conjecture for restricted elementary functions.
\textit{Annals of Mathematics}, 186(1), 237--275.
DOI: 10.4007/annals.2017.186.1.6.

\bibitem{BhattMorrowScholze2018}
B.~Bhatt, M.~Morrow \& P.~Scholze (2018).
Integral $p$-adic Hodge theory.
\textit{Publications Math\'ematiques de l'IH\'ES}, 128, 219--397.
DOI: 10.1007/s10240-019-00102-z.

\end{thebibliography}

\makeatletter
\let\test@bbl@sw\false@sw
\auto@bib@empty
\makeatother

\end{document}
